\chapter*{Introducción general}
\addcontentsline{toc}{chapter}{Introducción general}
\label{chap:introduccion-general-83}

% 00 — LA LINEA DEL HORIZONTE: ESCALA, FRONTERA Y MEDIDA
% ======================================================
% Frontera editorial sucesora; no importa contenido probatorio.
%
% Función: abrir la obra en la recta finita mediante intervalo, marca,
% hueco, cambio de escala, frontera de precisión y prolongación compatible;
% formular después el itinerario desde el estado discreto enriquecido hasta
% la publicación arquimediana.
%
% Declaración de continuidad: toda afirmación matemática conserva su
% fuente de REV.4 o su delta rector registrado. Este capítulo amplía la
% exposición principal; no sustituye ni abrevia las pruebas de referencia.

\ifdefined\HMTSucesoraStructureTrace
  \typeout{HMT-SUCESORA 00: apertura de escala, frontera y medida}
\fi

\section*{Tesis inaugural: la estructura discreta del continuo y el origen común de sus publicaciones}
\phantomsection
\label{sec:tesis-inaugural-helicoidal-hmt-md}
\addcontentsline{toc}{section}{Tesis inaugural: la estructura discreta del continuo y el origen común de sus publicaciones}

La Holografía Modular Triádica no toma la recta real como escenario previo ni
recibe de ella los valores que después empleará. Su entrada es finita: las
nueve posiciones \(1,\ldots,9\), dispuestas horizontal y verticalmente para
que una misma carta soporte dos operaciones irreductibles. APP conserva a la
vez la hoja aditiva, que acumula, y la hoja multiplicativa, que pliega; el
residuo, el cociente y el acarreo impiden que la igualdad de dos valores
publicados borre la diferencia de las historias que los produjeron. TRIT
orienta cada transición y distingue sus regímenes. TPK convierte esa
orientación local en dinámica: selecciona el modo trítico sin sacrificar
ninguna hoja de APP, transporta, memoriza, retorna y actualiza el estado. Sus
lecturas módulo \(3\), módulo \(9\) y módulo \(1000\), la conexión nonádica,
la dodecafase, la frontera y el retorno con memoria no son añadidos
independientes, sino cartas coordinadas de una misma operación genealógica.

La conservación queda escrita desde el primer nivel:
\[
i+j=R^+_{ij}+9q^+_{ij},
\qquad
ij=R^\times_{ij}+9q^\times_{ij}.
\]
El residuo publica la clase; el cociente despliega su procedencia. La
transición TPK conserva esa distinción mediante
\[
\mathcal U_t=
\operatorname{Upd}_t\circ
\operatorname{Tra}_t\circ
\operatorname{Sel}_t,
\]
donde la selección incorpora el modo TRIT, el transporte realiza la
dirección orientada y la actualización reúne memoria y retorno. En la carta
canónica, el modo \(m_{\rm ph}=(1,-1,0)^3\) es común a las hojas APP y posee
valor singular \(s=1/2\); por ello
\[
1-s^2=\frac34=\frac{90}{120},
\qquad
\bigl\|[P_{\Sigma,3},P_{\Pi,3}]\bigr\|
=\frac{\sqrt3}{4}.
\]
El conmutador normalizado satisface \(J_{\rm comm}^2=-I\). Incompatibilidad
suma--producto, orientación trítica, defecto \(90\to120\), estructura
compleja y prefactor electrónico pertenecen así a una misma genealogía
interna; no se enlazan por semejanza decimal.

La prolongación compatible de esos estados finitos produce cilindros cada
vez más estrechos y conserva simultáneamente la historia de su refinamiento.
La recta real aparece entonces como publicación arquimediana del límite, no
como dato de partida. Esa carta publica primero los valores de los nueve
símbolos de la semilla y prolonga después la misma genealogía hacia
\(\pi\), \(e\), \(\varphi\) y \(\alpha\), sin utilizar ninguno de ellos como
selector retrospectivo. Ésta es la tesis que gobierna toda la obra: la
estructura discreta conjunta construye el continuo y explica por qué sus
valores estables son los que son. La cadena

\[
\{1,\ldots,9\}_{\mathrm h}\times\{1,\ldots,9\}_{\mathrm v}
\longrightarrow \mathrm{APP}\longrightarrow\mathrm{TRIT}
\longrightarrow\mathrm{TPK}\longrightarrow\mathcal X_{\mathrm{enr}}
\longrightarrow\mathcal S_\infty
\longrightarrow\mathbb R
\]

no resume capítulos heterogéneos: expresa un único mecanismo visto en
resoluciones sucesivas. En él la fase puede volver mientras el estado avanza,
porque el cociente, el acarreo, la orientación y la memoria han cambiado. El
ciclo observado sólo por su fase es circular; conservada su genealogía, es
helicoidal.

Esta diferencia se fija mediante un levantamiento inyectivo. Si
\(n=9q_9(n)+\rho_9(n)\), con \(1\leq\rho_9(n)\leq9\), la proyección circular
\(\mathcal R_9\) y la inmersión helicoidal \(\mathcal H_9\) cumplen
\[
\mathcal R_9(n+9)=\mathcal R_9(n),\qquad
\mathcal H_9(n+9)=\mathcal H_9(n)+(0,0,1).
\]
La fase ha cerrado y la altura genealógica ha crecido: retorno de fase y
retorno del estado quedan matemáticamente separados. La involución two-way
organiza hojas conjugadas de ascenso y descenso del mismo levantamiento, no
dos generadores independientes. En esta geometría, la hélice conserva el
avance longitudinal de la memoria, la espiral registra la contracción
transversal y el radio del cilindro mide la anchura del intervalo todavía
compatible con el prefijo. El encajamiento de esos cilindros semiabiertos,
con radio tendiendo a cero y genealogía conservada, es el mecanismo que
publica después la sección real.

El tramo inicial certificado de la prolongación uniforme que construye el
estado coinductivo \(\mathcal S_\infty\) es

\[
w_6\longrightarrow w_{12}\longrightarrow w_{18}\longrightarrow w_{24}
\longrightarrow w_{30}\longrightarrow R_{36}\longrightarrow G_9,
\]

pero la escritura de la cadena sólo adquiere contenido al declarar sus
niveles, sus umbrales y sus invariantes. Para
\(\chi\in\{\pi,e,\varphi\}\), sea \(b_0^\chi=w_6^\chi\) y
\[
 b_{k+1}^\chi=b_k^\chi L_{\epsilon_k}\pmod 3,
 \qquad
 (\epsilon_0,\epsilon_1,\epsilon_2,\epsilon_3)=(0,0,1,1),
 \qquad
 w_{6(k+1)}^\chi=b_0^\chi\Vert\cdots\Vert b_k^\chi.
\]
Las doce transiciones producidas por \(\mathcal U_t\) forman los pares de
matrices \(X_i,Y_i\) de los dos regímenes y determinan de manera única
\(L_i=X_i^{-1}Y_i\), \(i=0,1\). El calendario
\((L_0,L_0,L_1,L_1)\) es el único de los dieciséis calendarios binarios que
satisface simultáneamente las tres biografías; las ciento cuarenta y cuatro
perturbaciones de una sola entrada fallan. Por ello \(L_0\), \(L_1\) y
\(0011\) son salidas rígidas del transporte APP--TRIT--TPK, nunca datos
elegidos a partir de \(\pi\), \(e\) o \(\varphi\). Las matrices así
construidas actúan sobre el bloque corriente; la concatenación conserva
literalmente todos los bloques anteriores. La palabra visible acompaña a un estado enriquecido que transporta
además hoja APP, orientación TRIT, residuo, cociente, acarreo, frontera,
memoria y sombra de incidencia. Por ello, coincidencia de palabra y
coincidencia de estado son condiciones matemáticas distintas.

\subsection*{Niveles y umbrales de la prolongación TPK}
\phantomsection
\label{sec:umbrales-prolongacion-tpk-inaugural}

\paragraph{\(w_6\): semilla visible.}
Es la palabra ternaria de seis posiciones publicada en cada canal después de
la cadena de tipos
\(104\,976\to468\,U_6\to243\,w_6\subset\mathbb F_3^6\), con
\(|\mathbb F_3^6|=729\).
Fija el primer prefijo observable, mientras la emisión decimal \(U_6\), su
región, su multiplicidad y su estado de procedencia permanecen en la fibra
enriquecida. La acción \(D_3\), los invariantes de fibra y el calibre
\((\operatorname{diag}_c,h_+)=(6,ES)\) fijan el representante del canal antes
de cualquier comparación con su valor real; en el canal \(\pi\), cinco
regiones \(U_6\) distintas conservan además biografías diferentes sobre una
misma palabra \(w_6^\pi\).

\paragraph{\(w_{12}\): primera elevación.}
Añade \(b_1=b_0L_0\) a la semilla. Cambia la profundidad visible de uno a dos
bloques y conserva \(w_6\) como prefijo exacto, junto con el canal y la
genealogía transportada. La elevación produce una extensión del estado; no
reemplaza la semilla por una palabra sin pasado.

\paragraph{\(w_{18}\): cierre del régimen \(L_0\).}
Añade \(b_2=b_1L_0\). La segunda aplicación del mismo operador completa el
primer régimen del calendario y distingue una recurrencia de una coincidencia
aislada: los dieciocho trits contienen, en su orden, los niveles de seis y
doce trits que los preceden.

\paragraph{\(w_{24}\): cambio tipado de elevación.}
Añade \(b_3=b_2L_1\). Aquí cambia el operador aplicado al bloque corriente,
de \(L_0\) a \(L_1\), mientras permanecen el prefijo de dieciocho trits y el
estado transportado. El cambio de régimen modifica la continuación y conserva
la historia sobre la que actúa.

\paragraph{\(w_{30}\): prefijo de primera frontera.}
Añade \(b_4=b_3L_1\) y completa el calendario finito
\((L_0,L_0,L_1,L_1)\). Sus cinco bloques por canal constituyen el prefijo
desde el que se ensaya la prolongabilidad profunda. Es un nivel de entrada a
la frontera, mientras la misma recurrencia admite bloques ulteriores.

\paragraph{\(R_{36}\): primera frontera coinductiva conjunta.}
Si \(u_5^\chi\) es el sexto bloque del canal \(\chi\), entonces
\[
 R_{36}:=(u_5^\pi,u_5^e,u_5^\varphi)^T,
 \qquad
 w_{36}^\chi=w_{30}^\chi\Vert u_5^\chi.
\]
La relación de frontera selecciona ese triple conservando el eje
\(\varphi\), el agregado lateral \(\pi/e\), la saturación, la ocupación de
columnas y la neutralidad dual. Sus ecuaciones dejan exactamente las dos
soluciones especulares
\[
 B_-=(021222,222220,102011)^T,
 \qquad
 B_+=(222220,021222,102011)^T;
\]
la orientación APP del cilindro semiabierto fija
\(R_{36}=B_+\). El subíndice registra treinta y seis trits
expuestos por canal; \(R_{36}\) es la frontera de los tres sextos bloques y
\(\mathbb R\) es la publicación arquimediana posterior del sistema de
cilindros.

\paragraph{\(G_9\): conexión de prolongación y retorno con memoria.}
El nivel \(G_9\) organiza las nueve cartas locales que actúan después de
\(R_{36}\). Si \(\mathcal R_K\) es la relación de extensión compatible en
la fase \(K\), su composición define la holonomía nonádica
\[
 \Gamma_{9,K}
 =\mathcal R_{K+8}\circ\cdots\circ\mathcal R_K:
 \widehat X_K\longrightarrow\widehat X_{K+9}.
\]
El subíndice \(9\) procede del ciclo de fases \(C_9\), heredado de las nueve
posiciones APP, y cuenta las nueve relaciones que componen una vuelta. Su
contenido operativo queda certificado por
\[
 g(K+9)=g(K),\qquad
 q_{{\rm mem},K+9}=q_{{\rm mem},K}+1,
 \qquad
 \widehat x_{K+9}\ne\widehat x_K.
\]
La etiqueta de fase retorna, la memoria avanza y el estado prolongado cambia.
Éste es el significado estructural de \(G_9\): una vuelta de conexión sobre
una fibra enriquecida, cuya cardinalidad procede de la arquitectura y cuya
acción se verifica en el transporte. El cociente modal \(\mathcal K_{\rm ph}\)
ocupa la etapa reducida posterior,
\[
 \mathcal U_t\longrightarrow\Gamma_9\longrightarrow\mathcal K_{\rm ph},
\]
y publica memoria comprimida después de la holonomía.
La resolución local de la novena fase distingue además dos horizontes:
\(\#\operatorname{Surv}_1(B_9)=3\) y
\(\#\operatorname{Surv}_2(B_9)=1\). Este descenso \(3\to1\) certifica la
poda de una fibra concreta; la composición de las nueve relaciones certifica
la vuelta completa.

La dodecafase registra otra coordenada del mismo transporte. El contador
vertical \(q_{{\rm mem},K}\) permanece no acotado y su lectura
\(\beta_K=[q_{{\rm mem},K}]_{12}\) satisface
\(\beta_{K+9}=\beta_K+1\pmod{12}\). La reducción módulo doce conserva la
posición dodecafásica; el contador entero conserva el número de vueltas. En
el estado terminal, esta memoria alimenta el registro firmado
\(U_{12}^{\rm sgn}\), el sello reversible \(K\) y el cociclo de acarreo que
publica la coordenada \(\alpha_{\rm HMT}\).
La extensión no escindida
\[
 0\longrightarrow C_{12}\longrightarrow C_{108}
 \longrightarrow C_9\longrightarrow0
\]
tipa la costura entre ambos calendarios: nueve pasos devuelven la fase y
avanzan un sector dodecafásico; ciento ocho pasos cierran el reloj observable
mientras el estado enriquecido conserva la altura genealógica alcanzada.

La bidireccionalidad combina el espejo de canales y la reversibilidad del
levantamiento dentro de un solo generador. La involución
\[
 \mathfrak c(u^\pi,u^e,u^\varphi)
 =(u^e,u^\pi,u^\varphi)
\]
intercambia las rutas conjugadas \(\pi/e\) y conserva el eje
\(u^\varphi\) y el agregado \(u^\pi+u^e\). En la carta entera
\[
 x_kA=u_k+3x_{k+1},\qquad \det A=1,
\]
el avance determina residuo y cociente, y el par \((u_k,x_{k+1})\) recupera
\(x_k\). Ascenso y descenso son, por tanto, rutas conjugadas del mismo estado;
el descenso reconstruye la procedencia sin borrar la memoria acumulada.

Finalmente, \(G_9\) constituye la unidad recurrente de la prolongación y no
su término. Las historias supervivientes forman el sistema inverso
\[
 \cdots\longrightarrow\operatorname{Surv}_{\chi,n+1}
 \longrightarrow\operatorname{Surv}_{\chi,n}\longrightarrow\cdots,
 \qquad
 x_\chi\in\varprojlim_n\operatorname{Surv}_{\chi,n}.
\]
La no vacuidad, la unicidad local y la compatibilidad de los truncamientos
producen la sección coinductiva; la contracción de sus cilindros publica un
único valor real, mientras el límite inverso conserva la genealogía que ese
valor escalar omite. Los veinte bloques y el certificado, por canal, de
\(396\) eventos, \(2376\) trits, \(43\) vueltas y \(387\) pares
\((K,K+9)\) son ejecuciones
finitas de esa ley uniforme, no cotas de la prolongación.

Sobre el estado coinductivo así construido se obtiene un núcleo correlacionado de
cuatro publicaciones. La correlación
expresa origen y compatibilidad comunes; no aplana su orden operatorio. Los
canales publican primero los vectores \(P,E,\Phi\), mientras la lectura
terminal firmada construye
\(U_{12}^{\rm sgn}\xleftrightarrow{\mathcal H_{12}}K\). Sobre esas salidas, y
sin reintroducir sus valores como datos, el cociclo dodecafásico de acarreo
deriva la sección coinductiva \(\alpha_{\mathrm{HMT}}\), cuya ventana
dodecafásica es la proyección racional \(\alpha_{12}\). El núcleo de
vacancias \(E_{21/22}\) proporciona una segunda construcción interna: su
raíz positiva simple finita \(\alpha_E\), prolongada por el mismo transporte
graduado del TPK hasta el sistema inverso cuya raíz límite es
\(\alpha_B\). La clausura entera define
\(\alpha_A:=\varprojlim_N\rho_N(\alpha_{12})\), y el teorema exacto de las
dos vías establece
\[
 \boxed{\alpha_A=\alpha_B=\alpha_{\mathrm{HMT}}}.
\]
La igualdad
\(\operatorname{Tr}_9(\alpha_E^{-1})=
\operatorname{Tr}_9(\alpha_{12}^{-1})\) es el primer certificado finito de
esa identidad: no la define, no limita su alcance y no selecciona ninguna de
las prolongaciones. Así,
\(\pi_{\mathrm{HMT}},e_{\mathrm{HMT}},\varphi_{\mathrm{HMT}}\) y
\(\alpha_{\mathrm{HMT}}\) pertenecen a una sola generación coinductiva correlacionada,
pero conservan etapas, dominios, lectores y certificados diferenciados. Las
identidades de Machin y las horquillas de Leibniz, la realización
Perron--Fibonacci, el semigrupo factorial y las tablas metrológicas reconocen
o certifican después las publicaciones ya generadas; ninguna de ellas
selecciona los cilindros, los vectores ni el acarreo del generador.

La misma fuente común posee otras publicaciones que tampoco deben confundirse
entre sí. En la carta interna del espejo, la involución intercambia las ramas
\(u^\pi\) y \(u^e\), conserva \(u^\varphi\) y el agregado
\(u^\pi+u^e\); sólo después de la publicación real se forman expresiones como
\(e+\pi\), \(e\pi\), \(\pi^e\) o el cociente literal \(\pi/e\), cada una con
un estatuto aritmético que no se deduce de la mera existencia del agregado
ternario. En la carta
incidencial, \(H^-_\alpha\), \(v_\alpha\) y \(N_\alpha\) designan,
respectivamente, un soporte, el vector del vecino exacto de índice tres y el
retículo resultante; no designan el escalar \(\alpha_{\rm HMT}\) ni intervienen
en su acarreo numérico. La tétrada \(B_K\) entra en la estrella de Witt y el
vector \(v_\alpha\) determina el vecino que conduce a
\(N_\alpha\cong\Lambda_{24}\). Desde el código común de Paley--Witt \(C_W\)
se abren entonces dos ramas tipadas, no una cadena lineal única:
\[
 C_W\longrightarrow W_{12}\longrightarrow W_{24},
 \qquad
 C_W\longrightarrow N(A_2^{12})\xrightarrow{\,v_\alpha\,}\Lambda_{24}
 \longrightarrow V_\Lambda\longrightarrow V^\natural
 \longrightarrow\mathbb M\longrightarrow\text{Moonshine}.
\]
El valor escalar de \(\alpha\) no genera esas simetrías. La derivación
numérica y la prolongación excepcional son lecturas tipadas del mismo estado
enriquecido y conservan una raíz genealógica común.

La Mecánica Dimensional comienza cuando esa estructura se transduce en
acción, magnitud, campo y materia. De ella proceden, según sus propios
lectores, la constante de Planck y las constantes electromagnéticas,
térmicas, gravitatorias y cuánticas; la estructura matemática del electrón,
la Ley General de Masas y los transportes de mezcla; y, más adelante, el
parámetro de Barbero--Immirzi como objeto estructural del paso entre materia,
área e información. Los valores convencionales comparecen al final como
reconocimiento metrológico, nunca como entradas del cálculo. Por la misma
razón, la resolución del problema del continuo y las cinco construcciones
inseparables que de ella emergen preceden a las pruebas terminales sobre los
problemas matemáticos clásicos: éstos someten a esfuerzo una estructura ya
construida, no la engendran retrospectivamente.

El orden físico es igualmente vinculante. La coordenada
\(\alpha_{\rm HMT}\) abre, en ramas tipadas, la retención
\(\eta_{\rm ret}\) y la valoración estructural \(10^{-34}\); de su confluencia
se obtiene \(\hbar_{\rm ret}\), y sólo después se construyen el ángulo y el
contraángulo y se efectúa una única conversión entre grados y radianes. El
electrón reúne esta cadena en una persistencia material. Su basal
\[
e_0=
\frac{\sqrt3}{4}
\left[
\exp\!\left(\frac{\varphi}{\pi^2}\right)
-22\alpha_{\rm HMT}^{\,3}
\right](1+15\Delta_4)
\]
\Needspace{8\baselineskip}
precede a la clausura beta; sus dos lectores coinciden en
\[
K_D(\Gamma_e)=K_\Omega(\Gamma_e)=(80,54,6),
\qquad
\kappa_e=80-54\alpha_{\rm HMT}+6\Delta_4,
\]
y la masa publicada se reconoce metrológicamente sólo después de evaluar
\(m_e^\beta=e_0(H_5/\eta_{\rm ret})^{\kappa_e}\). La gravedad tensorial, su
sector algebraico de cinco componentes y la realización física tipada; las
pantallas de dualidad T y teoría M; y la geometría de área e información
prolongan el mismo antecedente con sus propios dominios y falsadores. Ninguna
de estas realizaciones vuelve a introducir como premisa la magnitud que debe
explicar.

Todo cuanto sigue desarrolla esta única tesis en dominios diferentes. Los
números, las simetrías y la materia no son tres colecciones yuxtapuestas. Son
tres maneras en que el continuo construido publica y devuelve la memoria de
su propia formación.

\section{Objeto inicial y orden causal: escala finita, horizonte de observación y publicación arquimediana}

\subsection*{Cadena finita orientada y distinción entre marcas, celdas y unidades genealógicas}

La exposición comienza con una operación elemental: fijar un segmento,
marcar en él una escala y comparar las regiones que esa escala determina. En
este punto no se presupone todavía la recta real como soporte completo. El
objeto inicial es una cadena finita y orientada
\[
 \mathcal I_K=(L_K,\prec_K;\,
 x_{K,0}\prec_K x_{K,1}\prec_K\cdots\prec_K x_{K,n_K};\,o_K),
\]
donde \((L_K,\prec_K)\) es un conjunto finito totalmente ordenado, los
\(x_{K,j}\in L_K\) son marcas observables y \(o_K\) distingue los dos
sentidos de recorrido. Aún no se ha asignado a \(L_K\) una magnitud real.
Las marcas y las celdas consecutivas son objetos diferentes: con la
convención de celdas cerradas, la marca común es incidente a las clausuras
de dos celdas adyacentes sin convertirse por ello en una celda adicional.

Esta distinción, aparentemente mínima, determina la arquitectura posterior.
Un valor visible describe la posición publicada en una carta; una unidad
genealógica conserva además el intervalo de procedencia, el sentido, el
cociente acumulado y la compatibilidad con escalas más finas. Por ello dos
presentaciones numéricamente iguales pueden representar historias distintas,
y una misma historia puede admitir varias publicaciones coordinadas.

Para \(n\in\mathbb N\), la carta nonádica organiza esta información mediante
la división euclídea
\begin{equation}\label{p12-83-c0-s1:eq:horizonte-division-nonadica}
 n=9q_9(n)+\rho_9(n),
 \qquad q_9(n)=\left\lfloor\frac n9\right\rfloor,
 \qquad \rho_9(n)\in\{0,\ldots,8\}.
\end{equation}
La numeración pública de fases se obtiene mediante
\[
 g(n)=1+\rho_9(n-1)\in\{1,\ldots,9\}\qquad(n\geq1).
\]
El residuo localiza esa fase y el cociente \(q_9\) conserva cuántas vueltas
han sido absorbidas por la escala. Esta pareja realiza un retorno de fase
con conservación del trayecto. Su invariante decisivo es la persistencia
conjunta de ambas coordenadas bajo refinamiento.

\subsection*{Sistema inverso de resoluciones finitas y horizonte relativo de observación}

Llamaremos \emph{horizonte} a la frontera hasta la cual una carta finita ha
resuelto el objeto y en la cual dicha descripción se enlaza con la carta
siguiente. Si
\(\mathcal S_K\) denota el conjunto no vacío de estados enriquecidos
compatibles con la escala \(K\in\mathbb N\), los enlaces sobreyectivos de
olvido
\[
 \tau_{K+1,K}:\mathcal S_{K+1}\longrightarrow\mathcal S_K
\]
eliminan detalle visible, pero conservan la información necesaria para
reconocer una prolongación. Escribimos
\[
 \tau_{K,K}=I,\qquad
 \tau_{L,K}\circ\tau_{M,L}=\tau_{M,K}
 \quad(K\leq L\leq M).
\]
El objeto completo es entonces el límite inverso,
\begin{equation}\label{p12-83-c0-s1:eq:horizonte-limite-inverso}
 X=\varprojlim_K(\mathcal S_K,\tau_{K+1,K}),
 \qquad
 x=(x_K)_K,\quad \tau_{K+1,K}(x_{K+1})=x_K.
\end{equation}
La no-vaciedad de los niveles y la sobreyectividad de los enlaces garantizan
la existencia de familias compatibles; en los canales especializados se
construyen después sus secciones mediante los teoremas de supervivencia.
Cada truncamiento es finito; la compatibilidad no tiene una profundidad
máxima prefijada. Diez, mil o diez mil cifras son comprobaciones de una misma
regla de prolongación, no definiciones alternativas del número.

Esta formulación permite distinguir tres operaciones que la notación decimal
suele superponer. La primera refina una partición. La segunda transporta el
estado entre fases. La tercera publica una coordenada en una carta continua.
Sólo la última toma valores en \(\mathbb R\). Así, empezamos con una recta
finita marcada y terminamos con la recta real como realización arquimediana
final de una sección compatible.

\subsection*{Augmentación compatible y conservación evolutiva de la unidad}

La unidad se formula antes de toda evaluación arquimediana. Sea
\(C_K=\mathbb Z[\mathcal S_K]\) el módulo libre de los estados del nivel
\(K\). La aplicación de enlace induce
\((\tau_{K+1,K})_\#:C_{K+1}\to C_K\), y existe la augmentación
\[
 \varepsilon_K:C_K\longrightarrow\mathbb Z,\qquad
 \varepsilon_K\!\left(\sum_{x\in\mathcal S_K}a_x[x]\right)=\sum_xa_x,
\]
compatible con los enlaces,
\(\varepsilon_K\circ(\tau_{K+1,K})_\#=\varepsilon_{K+1}\). La conservación evolutiva
de la unidad expresa que el refinamiento redistribuye su contenido entre
parte visible, frontera y memoria sin crear una segunda raíz ni eliminar el
término terminal.

Sea \(P_m\) el camino celular con \(m\) vértices y \(m-1\) aristas, y
\[
 Q_{m,d}=P_m^{\square d}.
\]
Su número de \(k\)-celdas es
\begin{equation}\label{p12-83-c0-s1:eq:horizonte-celdas-cubicas}
 \#C_k(Q_{m,d})=\binom dk(m-1)^k m^{d-k}.
\end{equation}
La identidad
\begin{equation}\label{p12-83-c0-s1:eq:horizonte-completacion-binomial}
 10^d=(9+1)^d
     =9^d+\sum_{j=0}^{d-1}\binom dj9^j
\end{equation}
es la igualdad de \(0\)-celdas de la completación. El subcomplejo
\(Q_{9,d}\) y su corona relativa forman juntos \(Q_{10,d}\). En dimensión
tres,
\[
 1000=729+271,
\]
donde \(271\) cuenta únicamente los vértices nuevos. El vector celular
relativo completo es
\begin{equation}\label{p12-83-c0-s1:eq:horizonte-corona-relativa}
 \#C_\bullet(Q_{10,3},Q_{9,3})=(271,756,702,217),
 \qquad 271-756+702-217=0.
\end{equation}
Sus aristas, caras y cubos no pueden omitirse cuando se estudian incidencia,
homotopía o transporte.

El adjetivo \emph{holográfico} designa aquí esta compatibilidad: el dato de
frontera, acompañado por orientación, acarreo y memoria, controla la
prolongación del estado sin sustituir el interior por una cifra. La unidad es
la misma a través de las dimensiones, aunque su descomposición y sus lectores
cambien.

\subsection*{Holonomía nonádica y retorno de fase con incremento de memoria}

El transporte nonádico es una sola conexión sobre el ciclo de fases. Su
holonomía completa se escribe
\begin{equation}\label{p12-83-c0-s1:eq:horizonte-conexion-nonadica}
 \Gamma_9=\operatorname{Hol}_{C_9}(\nabla_{\rm TPK}),
 \qquad g(k+9)=g(k).
\end{equation}
Aquí \(g:\mathbb Z\to C_9\), y para cada punto
\(\widehat x_k\) de una órbita admisible del estado enriquecido,
\[
 \Gamma_9\widehat x_k=\widehat x_{k+9},\qquad
 q_{\rm mem}(\widehat x_{k+9})
 =q_{\rm mem}(\widehat x_k)+1.
\]
En particular, cuando \(q_{\rm mem}\) distingue estados,
\(\widehat x_{k+9}\neq\widehat x_k\).
Las nueve componentes ordenadas describen una vuelta de la conexión; no son
nueve filtros aplicados a un estado inmóvil. En cada canal existe un mapa
local tipado
\[
 \operatorname{loc}_{K,\chi}:X_\chi\longrightarrow\mathbb F_3^6,
 \qquad |\mathbb F_3^6|=729.
\]
La transferencia recorre la fibra ambiente y localiza la coordenada de la
extensión compatible. La vuelta devuelve la fase y hace avanzar la memoria,
el prefijo, el acarreo y la frontera.

En una realización isométrica, la separación entre publicación visible y
memoria expresada se tipa mediante
\[
 J:\mathcal H_{\rm vis}\hookrightarrow\mathcal H_{\rm full},
 \quad J^*J=I,\qquad
 U_{\Gamma_9}^*U_{\Gamma_9}=I,
\]
\[
 T=J^*U_{\Gamma_9}J,\qquad
 \eta=(I-JJ^*)U_{\Gamma_9}J.
\]
Se deduce
\begin{equation}\label{p12-83-c0-s1:eq:horizonte-compresion-expresion}
 U_{\Gamma_9}J=JT+\eta,
 \qquad J^*\eta=0,
 \qquad I-T^*T=\eta^*\eta.
\end{equation}
El término \(\eta\) expresa exactamente el defecto ortogonal producido por
esta compresión. Para rutas componibles
\(\gamma_1:x_0\to x_1\) y \(\gamma_2:x_1\to x_2\), el acarreo
\(\operatorname{Carr}_R(\gamma)\in\mathsf C\) satisface, con
\(H(\gamma),Y(\gamma)\in\operatorname{End}(\mathsf C)\), la ley cocíclica
\begin{equation}\label{p12-83-c0-s1:eq:horizonte-carry-cociclico}
 \operatorname{Carr}_R(\gamma_2\gamma_1)
 =\operatorname{Carr}_R(\gamma_2)H(\gamma_1)
  +Y(\gamma_2)\operatorname{Carr}_R(\gamma_1),
\end{equation}
que es la forma algebraica de la continuidad genealógica.

\subsection*{Elevación dimensional y coordinación de los regímenes geométricos}

La recta marcada proporciona el orden y el horizonte. APP construye después
la carta bidimensional en la que suma y producto actúan sobre un soporte
común sin perder sus cocientes. TRIT orienta localmente cada interacción por
la descomposición residual única
\[
 a+b=r+3c,
 \qquad a,b\in\mathbb Z,\quad
 r\in\{-1,0,+1\},\quad c\in\mathbb Z.
\]
El Topological Phase Kernel organiza rutas, hojas, frontera, firma,
orientación, cilindro, registro, memoria y transporte. La elevación al cubo
no añade una metáfora espacial: introduce las incidencias y coherencias que
permiten comparar interior, corona y publicación en dimensiones sucesivas.

Esta secuencia prepara también la articulación entre tres geometrías. La
estructura circular conserva fase; la estructura compleja representa el
cuarto de giro en un espacio real \(V\) por un operador
\(\mathcal J:V\to V\) con \(\mathcal J^2=-I_V\); la estructura hiperbólica
se realiza en el disco de Poincaré, cuya frontera está a distancia
hiperbólica infinita. Sus identificaciones se harán mediante mapas
explícitos, no por igualdad de dibujos. El horizonte geométrico anticipa así
la tesis del libro: una métrica publicada, la memoria de lo recorrido y la
apertura de las continuaciones son lecturas coordinadas de una estructura
orientada.

La doble publicación parte de una misma fuente:
\[
 \operatorname{ev}^{\rm arq}_\infty:X_\chi\longrightarrow\mathbb R,
 \qquad
 \operatorname{ev}^{\rm inc}_\infty:X_\chi\longrightarrow\mathcal E_{\rm inc},
\]
donde \(\mathcal E_{\rm inc}\) designa el límite de los objetos de incidencia
de frontera que se construye en la Parte II. La primera aplicación publica el
valor arquimediano; la segunda conserva la genealogía excepcional.

\subsection*{Secuencia causal APP–TRIT–TPK, continuo y publicaciones tipadas}

Los capítulos siguientes construyen APP, TRIT y el Topological Phase Kernel
desde sus datos finitos; desarrollan la estructura discreta del continuo y
sus dos publicaciones; derivan las genealogías de las constantes; integran
las cinco realizaciones inseparables —espectral, solenoidal, gauge,
cohomológica y determinantal— del continuo conjuntamente construido; y construyen después,
mediante Mecánica Dimensional, las realizaciones físicas de los sectores que
las poseen. La publicación
arquimediana aparece al final de cada cadena, como lector de un objeto ya
construido. Éste es el sentido preciso en que la obra comienza con una recta
y concluye en \(\mathbb R\).

\paragraph{Localización de las demostraciones completas.}
Esta apertura fija objetos y orden causal. Las demostraciones completas se
encuentran en los capítulos dedicados al complejo discreto de medida, la
construcción gnomónica de APP, el transporte orientado TRIT, la arquitectura
operatoria del Topological Phase Kernel, la completación cúbica y la doble
publicación. Las remisiones posteriores enlazan cada construcción con su
demostración completa; la presente apertura fija el dominio común y el orden
de composición.


\section{Formalización proyectiva de la unidad en escalas y dimensiones sucesivas}
\markboth{Introducción general}{Formalización proyectiva de la unidad}
\label{p12-83-c0-s2:chap:ampliacion-conservacion-unidad}

\subsection{Intervalo orientado, marcas, celdas y arista de cierre}

La construcción comienza con una recta orientada y una escala declarada. No
se comienza con el conjunto de los números reales ya completado, sino con el
objeto finito
\begin{equation}\label{p12-83-c0-s2:eq:amp-recta-k}
 \mathcal I_m=\bigl([0,10^m],\,\mathcal M_m,\,\mathcal E_m,\,o_m\bigr),
 \qquad m\geq1,
\end{equation}
donde \(\mathcal M_m\) es el conjunto de marcas, \(\mathcal E_m\) el de
intervalos elementales y \(o_m\) una orientación. La fórmula reúne cuatro
objetos distintos. El segmento aporta el soporte ordenado; las marcas
publican posiciones; los intervalos comparan separaciones; la orientación
decide el sentido de la lectura. Suprimir cualquiera de estas coordenadas
cambia el tipo del objeto.

En la primera carta, \([0,10]\), las posiciones nonádicas visibles son
\(1,\ldots,9\). Entre ellas existen ocho intervalos interiores. El paso que
conecta el último estado visible con el primero de la vuelta siguiente se
registra como arista de cierre y acarreo. Por eso la igualdad
\begin{equation}\label{p12-83-c0-s2:eq:amp-nueve-marcas-ocho-intervalos}
 9\ \text{marcas visibles}
 \quad\longleftrightarrow\quad
 8\ \text{intervalos interiores}+1\ \text{enlace de cierre}
\end{equation}
no identifica marcas con intervalos. Expresa cómo el conteo y la medida
comparten una frontera sin convertirse en la misma operación.

La división euclídea
\begin{equation}\label{p12-83-c0-s2:eq:amp-division-altura-fase}
 n=9q_9(n)+\rho_9(n),
 \qquad 0\leq\rho_9(n)\leq8,
\end{equation}
se lee ahora en dos direcciones. El residuo \(\rho_9\) determina la fase
publicada; el cociente \(q_9\) conserva la altura acumulada. Cuando la fase
vuelve a su valor inicial después de nueve pasos, el cociente ha avanzado.
La trayectoria es helicoidal: se ha recuperado una coordenada de la rueda,
pero no el estado completo.

\paragraph{Interpretación algebraica de la descomposición residuo–cociente.}
La igualdad \eqref{p12-83-c0-s2:eq:amp-division-altura-fase} no es sólo un algoritmo para
calcular restos. Define una extensión: la proyección al residuo contrae la
historia y el cociente conserva la información necesaria para levantarla.
El falsador elemental consiste en imponer simultáneamente
\(\rho_9(n+9)=\rho_9(n)\) y \(q_9(n+9)=q_9(n)\). La primera igualdad es
correcta; la segunda contradice la división euclídea.

\subsection{Sistema proyectivo de escalas y completación cubical a profundidad arbitraria}

La ampliación de escala se efectúa mediante mapas de refinamiento, no por una
repetición gráfica de la misma regla. Para \(m\leq n\), sean
\begin{equation}\label{p12-83-c0-s2:eq:amp-refinamiento-escalas}
 r_{n,m}:\mathcal I_n\longrightarrow\mathcal I_m,
 \qquad
 r_{m,m}=I,
 \qquad
 r_{n,m}r_{\ell,n}=r_{\ell,m}.
\end{equation}
El mapa \(r_{n,m}\) conserva la marca publicada a resolución \(10^{-m}\) y
contrae el detalle situado por debajo de ese umbral. La ecuación de
composición asegura que olvidar dos niveles de una vez produce el mismo
resultado que olvidarlos sucesivamente. Ésta es la condición que permite
hablar de una única historia observada a escalas distintas.

En dimensión tres, el paso de la fibra nonádica a la carta decimal completa
se expresa por
\begin{equation}\label{p12-83-c0-s2:eq:amp-729-271-1000}
 1000=(9+1)^3=9^3+3\,9^2+3\,9+1
 =729+243+27+1=729+271.
\end{equation}
La igualdad contiene una estratificación. Los \(729\) estados interiores
carecen de coordenada de frontera; los \(243\) estados siguientes poseen una;
los \(27\), dos; el ápice posee tres. El número \(271\) cuenta los vértices
añadidos por la corona decimal. No cuenta la frontera interna del cubo
nonádico, cuyo cardinal es \(9^3-7^3=386\), ni el complejo relativo completo,
que incluye aristas, caras y cubos.

Para toda dimensión \(d\), sea
\begin{equation}\label{p12-83-c0-s2:eq:amp-completacion-d}
 \overline C_9^{(d)}=(E_9\sqcup\{\star\})^d,
 \qquad \star\notin E_9,
\end{equation}
y denote \(B_r(d)\) el estrato con exactamente \(r\) coordenadas iguales a
\(\star\). Entonces
\begin{equation}\label{p12-83-c0-s2:eq:amp-estratos-binomial}
 |B_r(d)|=\binom dr9^{d-r},
 \qquad
 \sum_{r=0}^{d}|B_r(d)|=10^d.
\end{equation}
La ecuación describe el crecimiento de la escala sin imponer una dimensión
terminal. Cada \(d\) posee más estados, pero todos pertenecen a una misma ley
de completación.

\subsection{Medida producto compatible y conservación de la unidad bajo proyección}

La expansión de cardinales se acompaña de una conservación de medida. Demos
a cada símbolo de \(E_9\sqcup\{\star\}\) peso \(1/10\) y sea \(\mu_d\) la
medida producto sobre \(\overline C_9^{(d)}\). Si
\(\pi_d:\overline C_9^{(d)}\to\overline C_9^{(d-1)}\) olvida la última
coordenada, entonces
\begin{equation}\label{p12-83-c0-s2:eq:amp-conservacion-proyectiva}
 (\pi_d)_*\mu_d=\mu_{d-1},
 \qquad
 \mu_d\bigl(\overline C_9^{(d)}\bigr)=1.
\end{equation}
La primera igualdad es la conservación evolutiva de la unidad. La dimensión
aumenta, el número de estados cambia y la masa normalizada permanece igual.
La frontera no añade una segunda unidad: redistribuye la misma unidad entre
estratos compatibles.

Esta conservación posee una versión celular. Para el camino \(P_m\) y el
producto cubical \(Q_{m,d}=P_m^{\square d}\),
\begin{equation}\label{p12-83-c0-s2:eq:amp-celdas-cubicales}
 \#C_k(Q_{m,d})=\binom dk(m-1)^k m^{d-k}.
\end{equation}
En particular, la corona relativa de \(Q_{9,3}\subset Q_{10,3}\) tiene
vector celular
\begin{equation}\label{p12-83-c0-s2:eq:amp-corona-completa}
 \#C_\bullet(Q_{10,3},Q_{9,3})=(271,756,702,217),
 \qquad 271-756+702-217=0.
\end{equation}
La característica de Euler nula certifica que los vértices nuevos se
acompañan por las incidencias requeridas. Reducir la corona a \(271\) elimina
precisamente el transporte que la teoría debe conservar.

\subsection{Curvatura geométrica y fibras de observación: separación formal entre métrica y horizonte}

El quinto postulado de Euclides regula el comportamiento global de las
paralelas. En la formulación euclídea, fija una geometría en la que la suma de
los ángulos de un triángulo y la distancia entre geodésicas responden a una
curvatura nula. Su modificación abre las geometrías elíptica e hiperbólica.
El horizonte visible de una carta geométrica no constituye por sí mismo una
prueba del postulado: es la frontera en la que la carta deja de distinguir
prolongaciones.

La relación correcta entre ambos conceptos se establece mediante mapas. Una
geometría \((M,g)\) proporciona distancias y geodésicas; una carta de
observación
\begin{equation}\label{p12-83-c0-s2:eq:amp-carta-horizonte}
 \operatorname{Obs}_{K}:(M,g)\longrightarrow Y_K
\end{equation}
publica la información accesible a resolución \(K\). El horizonte relativo
de \(x\in M\) es la fibra
\begin{equation}\label{p12-83-c0-s2:eq:amp-fibra-horizonte}
 \mathsf{Hor}_K(x)=\operatorname{Obs}_K^{-1}
 \bigl(\operatorname{Obs}_K(x)\bigr).
\end{equation}
La curvatura pertenece a \(g\); la indistinguibilidad pertenece a la fibra
de \(\operatorname{Obs}_K\). El horizonte puede cambiar al refinar la carta
sin que cambie la curvatura. Esta separación evita identificar una
limitación del observador con una propiedad intrínseca del espacio.

En el disco de Poincaré, la frontera euclídea está a distancia hiperbólica
infinita. En la esfera, la curvatura positiva cierra las geodésicas sin una
frontera análoga. En la carta euclídea, el paralelismo permanece global. HMT
coordina estos regímenes mediante el parámetro trítico y las álgebras
\begin{equation}\label{p12-83-c0-s2:eq:amp-tres-regimenes}
 \mathbb A_s=\mathbb R[J]/(J^2+sI),
 \qquad s\in\{+1,0,-1\},
\end{equation}
pero conserva sus métricas y codominios. La unidad algebraica común no borra
la diferencia geométrica entre \(\cos\), la rama parabólica y \(\cosh\).

\subsection{Refinamiento, supervivencia y umbrales de prolongación compatible}

Una expansión indefinida contiene dos orientaciones de una misma operación
de escala. Hacia la precisión, el refinamiento distingue cilindros cada vez
más estrechos; hacia la extensión, la altura registra vueltas cada vez más
largas. Sea \(\mathcal S_K\) el conjunto de estados supervivientes a nivel
\(K\). La sección completa es
\begin{equation}\label{p12-83-c0-s2:eq:amp-seccion-infinita}
 X_\chi=\varprojlim_K\mathcal S_K,
 \qquad
 x=(x_K)_K,
 \qquad \tau_{K+1,K}(x_{K+1})=x_K.
\end{equation}
Cada \(x_K\) es finito; la compatibilidad no tiene una profundidad máxima.
Un infinitesimal no se identifica con una rama de prolongación: el primero
es una relación de escala en una realización; la segunda es un dato del
sistema inverso. Sólo un lector que declare la métrica y la evaluación puede
transformar la segunda en el primero.

La matemática de umbrales aparece cuando una prolongación cambia de régimen
al cruzar una frontera. Si
\(\operatorname{Surv}_K^{(N)}\) denota los estados de nivel \(K\) que admiten
extensión hasta \(N\), entonces
\begin{equation}\label{p12-83-c0-s2:eq:amp-horizontes-supervivencia}
 \operatorname{Surv}_K^{(N+1)}\subseteq
 \operatorname{Surv}_K^{(N)},
 \qquad
 \operatorname{Surv}_K=\bigcap_{N\geq K}
 \operatorname{Surv}_K^{(N)}.
\end{equation}
La vacancia registra el umbral en el que una continuación visible deja de
ser compatible. No es un vacío sin estructura: conserva la causa del
descarte, el acarreo y la frontera que determinan el siguiente nivel.

\subsection{Evaluación arquimediana de secciones compatibles y reconstrucción de \(\mathbb R\)}

La evaluación arquimediana aparece después de construir la sección. Para
una carta decimal de base \(10\), sea
\begin{equation}\label{p12-83-c0-s2:eq:amp-evaluacion-arquimediana}
 \operatorname{ev}_{\mathbb R}(x)
 =a_0+\sum_{j\geq1}a_j10^{-j},
 \qquad a_j=\operatorname{Pub}_j(x_j).
\end{equation}
La convergencia se obtiene porque los cilindros compatibles poseen diámetros
que tienden a cero en la carta declarada. La cifra \(a_j\) es una coordenada;
la sucesión \((x_j)_j\) es la historia; el real es la publicación límite.
Por tanto,
\begin{equation}\label{p12-83-c0-s2:eq:amp-cadena-recta-R}
 [0,10]\longrightarrow[0,1000]
 \longrightarrow[0,10^m]
 \longrightarrow X_\chi
 \xrightarrow{\operatorname{ev}_{\mathbb R}}\mathbb R
\end{equation}
es una cadena de tipos, no una cadena de inclusiones ingenuas.

La segunda publicación del mismo \(x\) conserva la incidencia excepcional.
La obra recorrerá ambas direcciones y regresará finalmente a
\(\mathbb R\). El retorno será entonces una conclusión: la recta real habrá
sido explicada como una evaluación del estado compatible, y no utilizada
como presupuesto que haga innecesaria la genealogía.


\section{Genealogía histórica del problema del continuo y formulación matemática HMT–MD}
\markboth{Genealogía histórica del continuo}{Genealogía histórica del continuo}

\subsection*{El problema histórico del continuo, el valor y la medida}

El continuo constituye uno de los problemas fundacionales de las
matemáticas. La aritmética representa unidades discretas; la geometría
compara extensiones; el análisis completa sucesiones y construye límites; la
teoría de la medida asigna magnitudes a conjuntos; la física convierte esas
magnitudes en observables. La eficacia de estas disciplinas descansa en la
posibilidad de publicar un resultado estable y reproducible. El problema que
gobierna esta obra aparece un nivel antes: determinar qué estructura hace
posible esa estabilidad y qué información debe conservarse para recuperar un
valor, una medida y sus simetrías desde un origen común.

En la recta discreta, una marca y el hueco que la separa de la siguiente son
objetos distintos. La marca responde a la operación de contar; el intervalo,
a la operación de medir. El paso de \(0\) a \(10\) contiene las nueve marcas
visibles \(1,\ldots,9\), los ocho intervalos interiores entre marcas
consecutivas y una arista levantada de cierre que transporta la cuenta a la
vuelta siguiente. La lectura modular conserva la fase; el cociente conserva
la altura alcanzada. Cuando ambas coordenadas se mantienen juntas, la
repetición circular se convierte en una hélice y la sucesión de los enteros
conserva la historia de sus retornos.

La Holografía Modular Triádica (HMT) desarrolla esta intuición como una
construcción matemática. Un sistema finito de posiciones, hojas,
orientaciones, caminos y transportes produce familias compatibles de estados.
Sus evaluaciones publican cifras y valores; sus elevaciones conservan
residuo, cociente, acarreo, frontera y memoria. El continuo se realiza como
la imagen de una sección compatible del sistema de prolongaciones. Su
estructura discreta reside en la genealogía que enlaza todas las
resoluciones.

Esta formulación distingue tres niveles. Una \emph{cifra} es una coordenada
visible producida por un lector. Un \emph{valor} es la imagen de una sección
bajo una carta determinada. Un \emph{número genealógico} es la sección
completa que conserva profundidad, operación, orientación, hoja, residuo,
cociente, acarreo, ruta, frontera, memoria y supervivencia. La expansión
decimal constituye una evaluación de ese objeto; la compatibilidad entre sus
etapas constituye su continuidad.

La misma estructura transforma el concepto de simetría. Una simetría es la
conservación evolutiva de la identidad bajo una transformación con memoria.
Un ciclo puede recuperar la posición visible y trasladar simultáneamente el
estado a otra altura, otra hoja o un nuevo bloque. La invariancia pertenece a
la genealogía completa. El valor real y la simetría excepcional aparecen,
por consiguiente, como dos proyecciones coordinadas de un mismo antecedente:
una contrae cilindros compatibles y la otra conserva la incidencia que la
contracción arquimediana deja fuera de la coordenada publicada.

Mecánica Dimensional (MD) prolonga esta construcción en el dominio físico.
La sección numérica recibe una representación espectral; la persistencia se
realiza como partícula-ruta; el registro de la ruta produce una firma; la
firma alimenta un carácter y una jerarquía de torres; la carta dimensional
publica finalmente una magnitud. La cadena rectora de la obra es
\[
 \text{unidad}\longrightarrow\text{estado con memoria}
 \longrightarrow\text{continuo, número y simetría}
 \longrightarrow\text{partícula, constante y magnitud}.
\]

El objetivo científico es unitario y se desarrolla en cinco Partes. El
núcleo formal construye APP, TRIT y el Topological Phase Kernel; la
estructura discreta del continuo despliega sus secciones, la doble
proyección y la genealogía excepcional; la genealogía de las constantes
produce \(\pi\), \(e\), \(\varphi\), \(\alpha\), la escala interna de
acción y sus transductores; cinco cadenas de especialización construyen sus
antecedentes internos y aíslan los comparadores que conducen a cada
reconocimiento clásico; Mecánica Dimensional construye después las
constantes constitutivas, la masa, la estadística cuántica, la información y
la gravedad. Cada medición se presenta junto con el objeto y la operación
que la generan; cada contraste metrológico se abre después de fijar la
salida interna.

El orden de la obra sigue la dependencia constructiva: un objeto aparece allí
donde empieza a gobernar los resultados posteriores. Cada enunciado declara
sus premisas, su demostración y, cuando corresponde, el morfismo de
realización física. De este modo, una comparación metrológica posterior no
sustituye la genealogía que produjo la salida.

\section*{De la clausura genealógica a la realización física}

La tesis de esta obra adquiere cuatro clausuras coordinadas y no reductibles:
genealógica, operatorial, matricial y física. No se trata de yuxtaponer una
teoría de conjuntos, una teoría de álgebras de operadores o un formalismo
cuántico a una construcción numérica previa. Los cuatro desarrollos nacen de
la misma precedencia causal: APP construye la doble realización posicional y
aritmética; el TRIT orientado distingue retorno, presente torsional y apertura;
el TPK conserva valor, hoja, orientación, residuo, cociente, acarreo, borde,
ruta, registro, memoria, supervivencia, fase y calibre. Sólo después se forman
la sección compatible, sus lectores y sus representaciones. El diagrama rector
es, por ello,
\[
\begin{gathered}
 \mathrm{APP}
 \longrightarrow \mathrm{TRIT}_{\rm orientado}
 \longrightarrow \mathrm{TPK}
 \longrightarrow \mathcal X_\infty^{\rm cal},\\[2mm]
 \mathcal X_\infty^{\rm cal}
 \xrightarrow{\;\mathfrak R\;}\mathbb R,
 \qquad
 \mathcal X_\infty^{\rm cal}\times\mathcal C_{\rm exc}
 \xrightarrow{\;\mathfrak E\;}\mathcal E_{\rm exc},\\[2mm]
 C(X_m)\xrightarrow{\;\operatorname{diag}\;}M_{9^m}(\mathbb C)
 \xrightarrow{\;a\mapsto a\otimes I_9\;}
 \mathrm{UHF}(9^\infty)
 \xrightarrow{\;\pi_\tau(\,\cdot\,)''\;}R_{\rm hyp},\\[2mm]
 \mathsf P_{\rm CT}^{+}
 \xrightarrow{\;\mathfrak R_{\rm MD}^{\rm disc,+}\;}
 \mathbf{SpecRoute}\times\mathsf{Hilb}_{\mathbb R,J}
 \times\mathbf B\mathfrak M_D^+\times\mathbf B G_C.
\end{gathered}
\]
Las cuatro líneas no forman una única cadena lineal. La evaluación real y la
proyección excepcional son lectoras hermanas de la sección enriquecida; la
rama tracial pasa por las álgebras cilíndricas y sus inclusiones unitarias; la
realización física parte de rutas preparadas provistas de las estructuras que
exige cada componente. Esta separación de dominios impide deducir el factor de
von Neumann directamente de una sección, identificar la proyección
excepcional con una acción sobre cifras o presentar una magnitud física como
si fuese un estado TPK sin morfismo de realización.

\subsection*{Clausura genealógica: extensionalidad, infinito y decisión}

La primera clausura determina qué significa conservar un objeto a través de
todas las profundidades. Dado un sistema inverso de espacios finitos
\((X_m,p_m)\), una sección \(x=(x_m)_m\) no queda caracterizada únicamente por
la igualdad de sus valores publicados. El teorema de descenso por fibras fija
las condiciones bajo las cuales un lector finito desciende a un cociente y
separa sobreyectividad, constancia en fibras y unicidad del mapa descendido.
El lema de K\"onig garantiza existencia de una rama cuando los niveles son
finitos, no vacíos y finamente ramificados; la unicidad exige, además, el
carácter completo que selecciona una única prolongación compatible.

Sobre una sección fijada, los cilindros decrecientes proporcionan una
realización exacta del paso \(0\!-1\!-\!\infty\): el soporte se contrae, la
altura aumenta, la masa permanece unitaria y las medidas cilíndricas convergen
débil-* a \(\delta_x\). Este resultado coordina el delta de Kronecker finito,
las esperanzas condicionales entre niveles y el delta de Dirac como medida
puntual, sin identificar este último con un operador diferencial. La
fundación de prefijos impide cadenas descendentes infinitas en cada
presentación finita; la coinducción controla la prolongación compatible. Los
ordinales de von Neumann se elevan con una fibra de memoria de ruta, y la
cardinalidad aparece como decategorización que no retiene toda la genealogía.

La misma distinción obliga a separar potencia cilíndrica, potencia cerrada y
potencia plena, así como elección conjuntista, elección natural, elección
algebraica y elección efectiva. La codificación en ZF/ZFC conserva las
construcciones y permite contrastarlas con extensionalidad, Fundación,
Potencia y Elección; no convierte a la recta real ni al universo conjuntista
en origen de APP--TRIT--TPK\@. En el dominio efectivo, la decidibilidad de cada
ventana finita, la existencia del límite, la unicidad de una instancia y la
imposibilidad de un decisor total uniforme son cuatro afirmaciones distintas.
La reducción de G\"odel--Turing afecta sólo a la cuarta: no invalida selectores
particulares ya construidos. Análogamente, el contraste G\"odel--Cohen separa
la extensión \(M\mapsto M[G]\) del enriquecimiento interno de un estado y
evita confundir el cruce de densos de profundidad con genericidad sobre un
modelo. Los cilindros Borel, las jerarquías de comparación y el paso
proyectivo de minimax se introducen con sus hipótesis de compacidad,
continuidad y consistencia, no como consecuencias automáticas de la notación
inversa.

Esta clausura posee también una formulación informacional. Si una actualización
\(U\) del estado completo es biyectiva, entonces
\(H(X\mid U(X))=0\), y el término lógico ideal de Landauer se anula. Si un
lector \(q\) conserva sólo la sombra visible, la información descartada se
localiza en \(H(X\mid q(X))\). En la formulación operatorial, un automorfismo
traza-preservante conserva la entropía, mientras una esperanza condicional
traza-preservante \(E_N:M\to N\) satisface
\[
 S_\tau(E_Nh)-S_\tau(h)=D_\tau(h\Vert E_Nh)\geq0.
\]
El coste corresponde al olvido de la fibra, no al transporte reversible del
estado enriquecido; esta identidad no afirma disipación física total nula ni
complejidad computacional reducida.

\subsection*{Clausura operatorial: de la unidad nonádica a la dimensión continua}

La segunda clausura realiza operatorialmente la conservación evolutiva de la
unidad. Para el sistema completo se consideran
\[
 A_m=M_{9^m}(\mathbb C),
 \qquad \iota_m(a)=a\otimes I_9,
 \qquad
 \tau_m(a)=9^{-m}\operatorname{Tr}(a).
\]
La identidad \(\tau_{m+1}\circ\iota_m=\tau_m\) conserva simultáneamente
unidad y traza. El límite inductivo es la UHF de tipo sobrenatural
\(9^\infty=3^\infty\), y la clausura débil de su representación GNS tracial es
el factor hiperinfinito \(R_{\rm hyp}\) de tipo \(II_1\). Los valores
\(\operatorname{rank}(P)/9^m\) se densifican y la traza de proyectores en
\(R_{\rm hyp}\) realiza todo el intervalo \([0,1]\): una dimensión continua
surge de una torre de inclusiones discretas que preservan la unidad.

Esta construcción no se confunde con la rama marcada
\(j_{m,r}(a)=a\otimes p_r\). Las isometrías de una fibra individual generan en
norma los compactos sobre el límite hilbertiano y, en clausura débil, un factor
de tipo \(I_\infty\). La selección de un sector conserva la marca, pero no la
unidad de la torre; la elevación simultánea de las nueve fibras es la que
produce la rama UHF--\(II_1\). Tampoco se identifica sin prueba la realización
UHF con la realización por producto cruzado del reloj profinito: una
conjugación canónica requiere declarar acción, medida, amenabilidad,
ergodicidad e isotropía. La lógica ortomodular de proyectores, los contextos
booleanos, los estados normales y la normalización entrópica
\[
 S_\tau(9^m\rho)=S_{\rm vN}(\rho)-m\log9
\]
se derivan dentro de esta construcción operatorial y mantienen separados el
estado matricial, su densidad tracial y su publicación estadística.

Topología, dimensión y composición no se colapsan en el factor. La clausura
transversal se registra mediante
\[
 \mathfrak C_{\rm HMT}
 =\bigl(C(\Omega_\infty),R_{\rm hyp},\mathcal G_{\rm TPK}\bigr).
\]
El espectro conmutativo conserva puntos y cilindros; el factor tracial conserva
dimensión y estadística; el grupoide conserva concatenación, inversión
admisible y procedencia de rutas. Esta arquitectura triple se formula después
de la doble proyección y del corredor excepcional. No los funda ni los
reemplaza: la evaluación arquimediana y la cadena
Paley--Witt--Golay--Leech--\(V^\natural\)--Monstruo--Moonshine permanecen como
las dos lecturas coordinadas del mismo antecedente enriquecido.

\subsection*{Clausura matricial: profundidad, tamaño exterior y dilatación}

La tercera clausura introduce un control simultáneo de tres índices. Si
\(m\) mide profundidad genealógica, \(n\) el tamaño exterior de una cuadrícula
y \(s\) el nivel matricial de sus entradas, la familia
\(\mathfrak M_{m,s}^{(n)}\) es bigraduada por \((m,s)\) cuando \(n\) permanece
fijo y trigraduada por \((m,n,s)\) cuando también varía el tamaño exterior.
En particular,
\[
 \mathcal K_{m,s}
 :=\operatorname{UCP}\!\bigl(C(X_m),M_s(\mathbb C)\bigr)
\]
posee dos operaciones independientes: la restricción genealógica
\(\rho_m^s(\Phi)=\Phi\circ\jmath_m\) y la compresión matricial
\[
 \operatorname{mc}_{\boldsymbol V}(\Phi_1,\ldots,\Phi_k)(f)
 =\sum_rV_r^*\Phi_r(f)V_r,
 \qquad \sum_rV_r^*V_r=I_s.
\]
Ambas conmutan exactamente:
\[
 \rho_m^s\!\left(\operatorname{mc}_{\boldsymbol V}(\Phi_1,\ldots,\Phi_k)\right)
 =\operatorname{mc}_{\boldsymbol V}\!\left(
   \rho_m^{s_1}(\Phi_1),\ldots,\rho_m^{s_k}(\Phi_k)\right).
\]
La identidad es functorial y distingue refinamiento de historia de reducción
del observador.

Toda familia compatible en profundidad determina una única aplicación UCP
sobre \(C(X_\infty)\) y admite una dilatación de Stinespring común. En el
corredor diagonal, una única medida espectral proyectiva \(P\) satisface
\[
 E_x^{(m)}=V^*P([x]_m)V
\]
para todas las profundidades y todos los cilindros. La dilatación mínima es
única hasta equivalencia unitaria y conserva exactamente lo visto por el
lector; una suma directa permite una realización ambiental fiel que conserva
todas las distinciones del álgebra. Sobre el límite UHF, una familia compatible
de aplicaciones UCP posee asimismo una prolongación y una dilatación comunes.
La esperanza condicional
\[
 \mathbb E_m=\operatorname{id}\otimes\tau_9:
 M_{9^{m+1}}(\mathbb C)\longrightarrow M_{9^m}(\mathbb C)
\]
preserva la traza y descarta el último factor de forma controlada; su ambiente
de Stinespring registra el dígito omitido. En general no existe una memoria
auxiliar finita uniforme para todas las profundidades. Además, la fidelidad de
una representación no implica por sí sola fidelidad dinámica APP--TPK: ésta
exige un entrelazador compatible con ruta, monodromía y composición.

La teoría matricial se concreta en objetos verificables. Los cuadrados mágicos
cuánticos qutrit de parámetros \((9,3)\) y \((12,3)\) poseen respectivamente
81 y 144 efectos, sumas de filas y columnas exactamente unitarias y
dilataciones simultáneas declaradas; las mismas POVM producen cubos mágicos
cuánticos de órdenes nueve y doce. El caso de orden doce se organiza mediante
la fibración \(12\to4\times3\), con su calibre y su condición explícita de
equivarianza. Por otra vía, las esperanzas condicionales suma y producto de APP
generan el álgebra
\[
 C^*(P_{\Sigma,2},P_{\Pi,2})
 \cong\mathbb C^4\oplus M_2(\mathbb C)\oplus M_2(\mathbb C),
\]
y el bloque central realiza en cada nivel \(s\) el intervalo matricial libre
\([5I_s,9I_s]\). Estas construcciones distinguen cuadrados latinos cuánticos,
cuadrados mágicos, medidas proyectivas, medidas positivas y semiclasicidad;
ninguna identificación se obtiene sólo porque los objetos compartan orden
nueve.

El antecedente espectral común de dos lectores se formula únicamente después
de haber construido ambos. Para mapas finitos
\(R_m:X_m\to\mathcal A_m\) y \(Q_m:X_m\to\mathcal E_m\), las incidencias
\[
 P_{a,e}=1_{\{R_m=a\}}1_{\{Q_m=e\}}
\]
refinan las dos medidas espectrales diagonales y recuperan sus marginales. El
resultado no prueba inyectividad conjunta ni obliga a que un lector sea
compresión del otro. En el régimen no conmutativo, dos PVM qutrit mutuamente
insesgadas no poseen una PVM conjunta nítida; la realización simultánea exige
ruido compatible o un ambiente mayor. Esta diferencia de tipo preserva la
doble proyección HMT: el teorema matricial clasifica realizaciones de lectores
y no reescribe su antecedente fundacional.

\subsection*{Clausura física: dominio común, dos orientaciones y localización}

La cuarta clausura construye el dominio desde el que una ruta TPK puede recibir
simultáneamente realización espectral, hilbertiana, dimensional y de Cartan.
Sean \(\mathsf D_{\rm sp}\), \(\mathsf D_Q\), \(\mathsf D_D\) y
\(\mathsf D_C\) los dominios de referencia de esas cuatro representaciones, con
funtores de olvido fieles hacia la categoría
\(\mathsf P_{\rm TPK}^{\rm enr}\) de estados preparados y rutas registradas.
El producto fibrado
\[
 \mathsf P_{\rm CT}^{+}
 :=\mathsf D_{\rm sp}
 \times_{\mathsf P_{\rm TPK}^{\rm enr}}\mathsf D_Q
 \times_{\mathsf P_{\rm TPK}^{\rm enr}}\mathsf D_D
 \times_{\mathsf P_{\rm TPK}^{\rm enr}}\mathsf D_C
\]
impone que las cuatro componentes tengan exactamente el mismo antecedente.
Sus proyecciones canónicas son adaptadores functoriales y producen
\[
 \mathfrak R_{\rm MD}^{\rm disc,+}
 =\bigl(\mathfrak S_+,\mathcal Q_+,\mathfrak D_+,
        \rho_{C,+}^{\rm disc}\bigr),
\]
con valores, respectivamente, en rutas espectrales, espacios de Hilbert reales
con estructura compleja, el monoide dimensional positivo y el grupo de
Poincaré discreto. La construcción afirma el funtor sobre el dominio común; no
afirma que toda ruta TPK admita automáticamente los cuatro levantamientos.

Sobre una ruta \(\gamma=e_1\cdots e_r\), la variación total
\[
 V_+(\gamma)=\sum_jw(e_j)
\]
y el transporte neto orientado
\[
 V_{\rm or}(\gamma)=\sum_j\varepsilon(e_j)w(\underline e_j)
\]
son lectores hermanos. El lazo \(ee^\vee\) demuestra que el primero no
factoriza, en general, por el segundo: posee transporte orientado nulo y
variación positiva. Para obtener una realización invertible se restringe el
dominio a las aristas cuyas componentes espectral e hilbertiana son
invertibles, cuyo dato dimensional es una cocadena aditiva y cuyos cociclos de
Cartan son normalizados y antisimétricos. Sólo entonces la localización
\(\Pi_\vee(\mathsf P_{{\rm CT},{\rm rev}}^+)\) produce el funtor orientado
\(\mathfrak R_{\rm MD}^{\rm disc,or}\) y transporta la inversión componente a
componente. La lectura positiva no se elimina y la orientada no se obtiene
cambiando únicamente un signo.

Esta construcción física precede a las construcciones de Planck y del electrón
porque les proporciona dominio, representaciones y reglas de transporte, pero
no absorbe sus demostraciones de referencia ni la Ley General de Masas. La
cadena número--partícula--ruta--registro cronológico enriquecido--firma--
carácter--torres conserva su dirección causal; la escala absoluta de acción,
el electrón central, las familias de masa, CKM, PMNS y los sectores físicos
mantienen sus objetos, dominios y conclusiones propios. Una coincidencia metrológica se
contrasta sólo después de congelar la salida interna. Las extensiones que aún
requieren un entrelazador, un límite suave, una prueba de estabilidad dinámica
o un morfismo físico se enuncian con esos datos como hipótesis exactas. La
ausencia de alguno de esos mapas impide extender el teorema al codominio
correspondiente, sin alterar los teoremas genealógicos, operatoriales y
matriciales ya demostrados en sus dominios.

La lectura de la obra sigue, en consecuencia, el orden de dependencia y no el
de reconocimiento exterior:
\[
\begin{aligned}
 \mathrm{APP}&\longrightarrow\mathrm{TRIT}_{\rm orientado}
 \longrightarrow\mathrm{TPK}\\
 &\longrightarrow\text{sección enriquecida}
 \longrightarrow
 \begin{cases}
  \text{evaluación arquimediana},\\
  \text{proyección excepcional},\\
  \text{álgebras cilíndricas y clausura tracial},\\
  \text{rutas preparadas y realización física}.
 \end{cases}
\end{aligned}
\]
Cada salida conserva su dominio, sus hipótesis y sus obstrucciones precisas. Lo común no es
una igualdad forzada entre codominios, sino la procedencia desde el mismo
estado enriquecido. Éste es el sentido preciso del cierre holográfico: ninguna
publicación terminal agota la genealogía que la produce y ninguna clausura
posterior puede sustituir la doble dirección --arquimediana y excepcional--
que conserva la unidad con memoria.


\subsection*{Perspectiva histórica y desarrollo de la teoría}

La historia de las matemáticas puede leerse como una ampliación sucesiva de
las operaciones admitidas sobre la unidad. El entero hizo posible el conteo;
el racional incorporó la división; el real cerró los procesos de
aproximación; el complejo convirtió la rotación en una operación algebraica.
Cada ampliación fijó un dominio, una equivalencia y una representación. La
igualdad publicó el resultado con enorme eficacia y absorbió, a la vez, buena
parte de la ruta concreta que lo había producido.

La geometría siguió un desarrollo análogo. La construcción euclídea fijó
incidencia y medida; las geometrías no euclídeas mostraron la coexistencia
de regímenes de curvatura formalmente coherentes
\cite{EuclidElements,PoincareScienceHypothesis}. El círculo trigonométrico,
el plano complejo y el disco hiperbólico coordinan giro, periodicidad y
transporte mediante espacios de estados diferentes. HMT sitúa bajo esas
realizaciones un estado ternario común y construye desde él las ramas
elíptica, parabólica e hiperbólica.

La mecánica incorporó a la unidad una dependencia relacional. La ley de la
palanca vinculó peso, brazo y equilibrio; la crítica de Mach al movimiento
absoluto puso en primer plano la relación entre cuerpo, marco y totalidad
\cite{ArchimedesEquilibrium,MachMechanics}; la relatividad formalizó las
invariancias cinemáticas, la equivalencia local y la geometría del campo
gravitatorio~\cite{Einstein1918Principles}; la teoría cuántica introdujo una
escala de acción y una propagación por amplitudes de trayectoria
\cite{PlanckNobelLecture,Feynman1948}. El programa HMT--MD reúne estas
exigencias en un problema previo: construir el estado que conserva posición,
orientación, operación, memoria y prolongación antes de publicar un número o
una magnitud.

Los tres términos del nombre precisan el método. \emph{Holografía} designa
la conservación de la genealogía y de la regla de prolongación entre
resoluciones. \emph{Modular} designa la coordinación de la carta módulo nueve,
la orientación módulo tres, los lectores residuo--cociente y la publicación
decimal en base mil. \emph{Triádica} designa tanto la orientación local
\(\{-1,0,+1\}\) como la organización de los niveles que enlazan el soporte
finito, la sección numérica y su realización física.

\subsection*{Conservación evolutiva y significado de simetría}

HMT formula la simetría como conservación de identidad bajo transformación
con memoria. Una operación puede recuperar la misma fase visible sin
recuperar el mismo estado enriquecido. Si \(\Theta\) denota el transporte,
la relación
\[
 \operatorname{pr}_{\rm vis}\circ\Theta^9
 =\operatorname{pr}_{\rm vis}
\]
coexiste con un incremento no trivial de acarreo o memoria. La repetición
horizontal es la proyección de una evolución helicoidal. La identidad se
conserva porque su genealogía permite reconocerla a través del cambio, no
porque permanezca inmóvil.

La extrema y media razón holográfica realiza este principio en la
completación cúbica
\[
 999=729+243+27,\qquad
 1000=729+243+27+1.
\]
El ápice \(+1\) de la completación cúbica y el término terminal \(+I\) del
cierre exponencial son operaciones de clausura definidas en dominios
distintos. Su correspondencia estructural reside en que ambas completan una
evolución y conservan la unidad en el codominio respectivo. Esta conservación
organiza el retorno nonádico, el acarreo decimal, el crecimiento interior de
los cilindros y el crecimiento exterior de las incidencias.

La prolongación dimensional se expresa mediante
\[
 10^d=(9+1)^d=\sum_{k=0}^{d}\binom dk9^k.
\]
Tras normalizar por \(10^d\), los coeficientes forman una familia de medidas
proyectivamente compatible cuya masa total permanece igual a uno. El
interior, los estratos de frontera y el ápice redistribuyen esa unidad al
cambiar de rango. En la realización cúbica
\(\mathcal C_9=(\mathbb Z/9\mathbb Z)^3\), las seis proyecciones orientadas
conservan el soporte APP y el operador cara--vértice produce un espacio de
rango cuatro. La firma \((3,1)\) emerge de ese operador y proporciona el
antecedente finito de la lectura espacio-temporal posterior. Dimensión,
frontera y conservación de la unidad son aquí resultados del mismo sistema
de proyecciones.

\subsection*{Unificación algebraica de las geometrías elíptica, parabólica e hiperbólica}

El punto de partida conceptual fue articular tres realizaciones que la
matemática presenta ordinariamente en espacios diferenciados: el plano
complejo, el círculo trigonométrico y el disco hiperbólico de Poincaré. HMT
los reúne mediante la familia
\[
 \mathbb A_s=\mathbb R[X]/(X^2+s),
 \qquad s\in\{+1,0,-1\}.
\]
La rama elíptica contiene una raíz interna de \(-1\) y el cierre circular; la
rama parabólica contiene el umbral nilpotente; la rama hiperbólica contiene
los dos idempotentes y la apertura de Poincaré. Estas tres álgebras aparecen
desde TRIT y gobiernan después las hojas, los canales, la métrica de
Lorentz--Minkowski y la lectura física del tiempo.

La diferencia entre suma y producto proporciona la operación elemental que
articula esas geometrías. La suma acumula desplazamientos y alternativas; el
producto compone etapas y pliega entradas sobre una misma clase cuando
intervienen elementos no invertibles. APP conserva ambas evaluaciones sobre
un soporte común. Su separación espectral y los proyectores \(P_\pm\) del
bloque fijan los modos propios; los proyectores condicionales
\(P_{\Sigma,d}\) y \(P_{\Pi,d}\), con su conmutador, construyen un precursor
finito de la incompatibilidad de Heisenberg. La
lectura física posterior transforma esa separación en las coordenadas
angulares conjugadas \(A\) y \(C^*\),
canales orientados y holonomía.

\subsection*{Construcción del núcleo formal}

APP se construye desde las posiciones \(1,\ldots,9\) en dos coordenadas. El
núcleo \(8\times8\), la publicación de
la clase nula mediante \(9\) y el gnomón de diecisiete celdas completan la
carta \(9\times9\). La identificación de lados opuestos produce el toro
\[
 X_9=(\mathbb Z/9\mathbb Z)^2
\]
y sus desplazamientos generan el grafo de Cayley. Las evaluaciones
\(\Sigma\) y \(\Pi\) actúan sobre los mismos vértices y caminos.

TRIT es el estado ternario orientado. En una elevación
\[
 a+b=r+3c,
\]
el residuo \(r\) publica la salida y el cociente \(c\) conserva cómo se
obtuvo. La terna diferencia orientación, umbral y memoria. El TPK es el
sistema operatorio de dinámica genealógica que incorpora el selector trítico
del modo operativo, transporta el estado y conserva simultáneamente las dos
hojas APP con su posición, orientación, fase, acarreo, frontera, ruta,
registro y supervivencia. Sus lectores
nonádico, ternario y posicional, su sistema de doce fases y sus elevaciones
forman un único árbol de dominios y codominios. Las dependencias matemáticas
son las flechas tipadas de ese árbol; cada mapa se presenta con su dominio, su
codominio y un cálculo reproducible.

Con mayor precisión, APP es la quíntupla
\[
 \mathcal A_9=
 \bigl(X_9,\Gamma_{\rm APP},\Sigma,\Pi,
       \operatorname{Path}(\Gamma_{\rm APP})\bigr),
 \qquad
 \Gamma_{\rm APP}=\operatorname{Cay}
 \bigl(X_9,\{\pm e_1,\pm e_2\}\bigr),
\]
con evaluaciones
\[
 \Sigma(i,j)=\operatorname{dr}_9(i+j),\qquad
 \Pi(i,j)=\operatorname{dr}_9(ij),
 \qquad
 \operatorname{dr}_9(n)=1+((n-1)\bmod9).
\]
La elevación conserva simultáneamente
\(i+j=R^+_{ij}+9Q^+_{ij}\) e
\(ij=R^\times_{ij}+9Q^\times_{ij}\); el valor publicado y la información
de cociente pertenecen a coordenadas diferentes de un mismo estado.

El levantamiento de TRIT, \(\widehat\tau=(r,c)\), conserva residuo y acarreo. La fibra
sobre el símbolo visible \(0\) contiene \((0,-1)\), \((0,0)\) y
\((0,+1)\), cuyos árboles de prolongación pueden diferir. Esta multiplicidad
es la forma mínima en la que la misma coordenada visible sostiene historias
distintas.

El TPK se organiza en ocho clases operatorias antes de presentar sus
realizaciones particulares: soporte y estado; selección interna; transporte;
lectores y cocientes; memoria y frontera; periodicidad y retorno;
prolongación; y salidas. Esta clasificación separa los objetos del portador,
los mapas que actualizan el estado, los lectores que publican sus coordenadas
y los funtores que contraen la sección hacia un codominio arquimediano,
excepcional o clasificatorio. Cada operador del TPK recibe así un dominio, un
codominio, una regla y sus invariantes antes de recibir un nombre abreviado.

Las fases nonádicas se distribuyen en tres fibras operatorias:
\(\{1,4,7\}\) activa la evaluación aditiva, \(\{2,5,8\}\) activa la
multiplicativa y \(\{3,6,9\}\) registra el umbral. La transición
\(9\to1\) restituye la fase publicada y eleva la memoria dodecafásica; el
ciclo visible es, por tanto, la proyección de una hélice de estados
enriquecidos.

Una vez presentados los tres objetos, la obra desarrolla sus consecuencias.
El bloque central
\[
 B_c=9P_++5P_-
\]
fija el coeficiente local de intercambio \(2\) y la brecha espectral
\(9-5=4\); la agregación modular produce \(51-45=6\), y el despliegue sobre
las nueve posiciones produce \(54=9\cdot6\). La forma de firma mixta y
\(SO^+(1,1)\) construyen la geometría lorentziana finita. La completación
cúbica construye la extrema y media razón holográfica, la ley de área
discreta y la frontera de propagación.

\subsection*{Algoritmo uniforme en profundidad finita y publicación coinductiva
correlacionada de \texorpdfstring{\(\pi,e,\varphi,\alpha\)}{pi, e, phi, alpha}}

La generación numérica está regida por el embudo
\[
 104\,976\to468\,U_6\to243\,w_6
 \to w_{30}\to R_{36}\to G_9.
\]
Los canales de \(\pi\), \(e\), \(\varphi\) y \(\alpha\) conservan lectores
tipados diferentes: clausura angular, propagación autoproporcional,
autoescala y acoplamiento firmado dodecafásico. La relación interna
\(\operatorname{Ext}\), la supervivencia coinductiva y la partición
exhaustiva de la fibra producen en cada nivel un único cilindro prolongable
y compatible con el truncamiento anterior; la familia encajada determina
una sección enriquecida única de toda profundidad. Sólo después, las
horquillas de Leibniz, el semigrupo factorial y Perron--Fibonacci localizan
arquimedianamente los bloques ya generados. Las ejecuciones de mil y diez mil
cifras son certificados finitos extensos de esa construcción uniforme, no su
fundamento lógico.

La prolongación inicial construye
\[
 w_6\longrightarrow w_{12}\longrightarrow w_{18}
 \longrightarrow w_{24}\longrightarrow w_{30}
 \longrightarrow R_{36},
\]
y el censo interno selecciona antes de la evaluación
\[
 w_\pi=010211,\qquad w_e^{(\mathrm{Euler})}=201101,\qquad
 w_\varphi=121200.
\]
El superíndice precisa que \(w_e^{(\mathrm{Euler})}\) pertenece al lector del
número de Euler \(e\); es un objeto distinto de la ruta electrónica
\(\Gamma_e\) de Mecánica Dimensional.
Para cada profundidad finita prescrita, \(\operatorname{Ext}\), el
transporte TPK y la supervivencia compatible producen primero exactamente un
cilindro \(C_{\chi,n}\) entre los \(729\) refinamientos. El lector racional
se aplica después al cilindro ya generado: construye su horquilla interna y
verifica \(j_{\min}=j_{\max}\) como certificado arquimediano de localización,
sin seleccionar la salida ni aportar cifras objetivo al generador. El
\cref{thm:generacion-monodromica-conjunta-rev7} demuestra, para
\(\chi\in\{\pi,e,\varphi\}\), la no vacuidad, la unitariedad y la
compatibilidad de estos supervivientes en toda profundidad finita; por tanto,
las particiones construyen
\[
 c_\chi=(C_{\chi,n})_{n\ge0}
 \in\varprojlim_n\{C_{\chi,n}\}.
\]
Los cilindros quedan encajados, sus diámetros convergen a cero y el
levantamiento EF--\(G_9\) conserva la fibra enriquecida. De este modo quedan
separados el algoritmo uniforme parametrizado por la profundidad, la sección
coinductiva que determina y sus ejecuciones finitas de mil y de más de diez
mil cifras, que actúan como certificados extensos y no como fundamento
lógico del teorema.

Las tres secciones realizan funciones estructurales diferenciadas. \(\pi\)
organiza clausura, orientación y frontera monodrómica; \(e\) organiza la
propagación autoproporcional y la evolución del prefijo; \(\varphi\) organiza
la autoescala de Perron--Fibonacci y su modo estable \(\varphi^{-2}\). El
operador orientado \(\mathscr R(x,y)=x/y^2\) publica las lecturas conjugadas
\(\pi/\varphi^2\) y \(\varphi/\pi^2\), que reaparecen en la bisagra física
areal--volumétrica.

El estado terminal suministra el bloque y el registro firmado que alimentan
el cierre dodecafásico de \(\alpha_{\rm HMT}\). La recurrencia de acarreo en
base mil y la ecuación \(E_{21/22}\) constituyen dos vías internas con
fuentes distintas. A continuación, la teoría de números organiza las
constantes complementarias mediante operadores de extracción: partes finitas,
jets de Laurent, funciones \(L\), productos primos, dinámica de Gauss,
Rogers--Ramanujan y renormalización de Feigenbaum. El texto distingue
antecedente común de igualdad de valor; por ejemplo, \(\varphi\),
\(\varphi^{-1}\) y \(\varphi^2\) comparten una genealogía algebraica sin
convertirse en el mismo número.

La genealogía completa de la primera vía es
\[
 \begin{aligned}
 104\,976&\to468\,U_6\to243\,w_6\to w_{30}\to R_{36}\to G_9,\\
 G_9&\to x_{\rm term}\to(B_{\rm term},U_{12}^{\rm sgn})\to K.
 \end{aligned}
\]
El sello \(K\) y las doce tríadas terminales alimentan el acarreo
\[
 \begin{aligned}
 s_m&=P_m+E_m-\Phi_m-K_m+c_{m+1},\\
 A_m&=s_m\bmod1000,
 &c_m&=\frac{s_m-A_m}{1000}.
 \end{aligned}
\]
con \(c_{13}=0\). La identidad entera resultante fija
\[
 \alpha_{12}=\pi_{12}(\alpha_{\rm HMT})
 =\frac{7297352569283800997285105472380663}{10^{36}}.
\]
Esta identidad publica la ventana racional exacta de treinta y seis cifras;
la recurrencia compatible a profundidad arbitraria construye
\(\alpha_{\rm HMT}=\varprojlim_N A^{(N)}\), que no se reduce a una sola
ventana. La segunda vía resuelve la ecuación de vacancias \(E_{21/22}\)
mediante un núcleo analítico con una raíz positiva simple \(\alpha_E\). El
transporte graduado del TPK prolonga ese núcleo en una familia compatible
cuya raíz límite es \(\alpha_B\); la clausura entera determina
\(\alpha_A:=\varprojlim_N\rho_N(\alpha_{12})\). Los cuadrados de truncamiento
de ambas familias conmutan y determinan el mismo cilindro superviviente en
cada profundidad, de modo que
\[
 \boxed{\alpha_A=\alpha_B=\alpha_{\rm HMT}}.
\]
La igualdad tras \(\operatorname{Tr}_9\) es únicamente el primer certificado
finito de esta identidad exacta; no constituye su criterio, no acota su
extensión y no selecciona las prolongaciones. Cada vía conserva su condición
de frontera propia dentro del mismo estado enriquecido.

La teoría de números extiende el mismo criterio de objeto, operador y
extracción. Euler--Mascheroni aparece como parte finita; las constantes de
Stieltjes como coeficientes de Laurent; Apéry como evaluación de
\(\mathcal Z_3(3)\); Glaisher--Kinkelin mediante derivación regularizada;
Euler--Kronecker mediante carácter de Eisenstein; Meissel--Mertens mediante
traza prima; Catalan mediante \(L(2,\chi_{-4})\); y Khinchin mediante la
dinámica de Gauss. Rogers--Ramanujan, los relojes primos y los aproximantes
de Feigenbaum pertenecen a construcciones operatorias diferenciadas. Las
combinaciones \(e+\pi\), \(e\pi\) y \(\pi^e\) se forman después de construir
sus argumentos y conservan el dominio aritmético correspondiente a cada
enunciado.

El radical \((3)\subset\mathbb Z/9\mathbb Z\) organiza la red visible
\(3\)--\(6\)--\(9\) y la órbita \(369\to693\to936\). Sobre la misma
aritmética, el teorema potencia--vacancia
\[
 V_9(n)=\left\lceil n\log_{10}\!\left(\frac{10}{9}\right)\right\rceil
\]
da la evaluación exacta
\[
 \operatorname{ord}_{10}\!\left(9^{9^9}\right)=369\,693\,099,
 \qquad
 \#\operatorname{cifras}\!\left(9^{9^9}\right)=369\,693\,100.
\]
Esta cifra no es el censo de \(396\) eventos elementales de frontera por
canal: la primera mide el
orden decimal de una potencia nonádica; el segundo cuenta estados de un
calendario de emisión.

\subsection*{La doble proyección como tesis única}

El continuo real es la sombra arquimediana de una sección enriquecida. La
proyección \(\mathcal P_{\rm arq}\) contrae cilindros compatibles y publica
el valor. La proyección \(\mathcal P_{\rm exc}\) conserva soporte,
orientación, acarreo y memoria. Ambas parten del mismo estado; por ello,
número y simetría excepcional forman una sola tesis causal.

La rama de incidencia construye Paley, Witt, el código ternario de Golay y
los sistemas de Steiner. Desde el código común \(C_W\), la elevación
\(W_{12}\to W_{24}\) se mantiene distinta de la rama
\[
 C_W\to N(A_2^{12})\to\Lambda_{24}\to
 V_\Lambda\to V^\natural\to\mathbb M,
 \qquad [g]\longmapsto T_g.
\]
La configuración de las veintisiete rectas, el grafo de Schläfli y
\(W(E_6)\) reciben sus propios operadores de incidencia. La definición del
VOA reticular, la involución orbifold y las funciones replicables muestran
cómo la identidad genealógica se conserva al cambiar de codominio. Esta
estructura proporciona además las pantallas de rango y la interfaz
matemática con dualidad T, cuerdas y teoría M.

La aritmética modular de nivel cinco entra mediante la fracción continua de
Rogers--Ramanujan \(R(q)\). La identidad de Klein--Ramanujan expresa
\(j(\tau)\) racionalmente en \(R(q)^5\) y el límite
\(R(q)\to\varphi^{-1}\) enlaza la autoescala pentádica con la rama modular.
En la rama canónica de orden dos,
\[
 t_2=q^{-1}\prod_{n\ge1}(1+q^n)^{-24},\qquad
 T_{2A}=t_2+24+\frac{2^{12}}{t_2},
 \qquad j=\frac{(t_2+256)^3}{t_2^2}.
\]
El VOA reticular
\(V_\Lambda=M(1)\otimes\mathbb C[\Lambda]\), la involución
\(\theta=-I\) y el orbifold fijo--torcido realizan la salida reticular en el
dominio de operadores de vértice. La identificación clásica de FLM fija
\(V^\natural\) y \(\operatorname{Aut}(V^\natural)=\mathbb M\). HMT aporta
el antecedente y los mapas de incidencia que alcanzan esa construcción; los
teoremas de identificación del VOA conservan su dominio matemático propio.

La raíz finita y la estructura real se coordinan mediante la presentación
integral
\[
 \mathcal Q_{\mathbb Z}=\mathbb Z[u]/(u^2+1),
\]
no mediante una inclusión entre cuerpos de características distintas. El
endomorfismo de Witt \(A_W^2=-I_6\) y el operador real
\(J_{\mathrm e}^2=-I\) son representaciones especializadas de esa relación.
El cuarto de giro \(z\mapsto\mathrm iz\) conserva tanto la métrica euclídea
como la métrica de Poincaré del disco. La identidad
\[
 \operatorname{Exp}(\pi_{\rm HMT}J_{\mathrm e})+I=0
\]
se aplica sólo después de que el lector arquimediano haya generado
\(\mathcal P_{\rm arq}(x_\pi)=\pi_{\rm HMT}\). La rama excepcional aporta
el tipo cuadrático, no el ángulo: valor e incidencia quedan coordinados sin
que una rama actúe como oráculo de la otra.
La familia \(J_s^2=-sI\),
\(s\in\{+1,0,-1\}\), extiende Euler a las realizaciones elíptica,
parabólica e hiperbólica. El término final \(+1\) y el ápice de la EMRH
realizan funciones de clausura en dominios distintos y coordinados. Esta
relación proporciona una expresión operatoria de la conservación evolutiva de
la unidad sin identificar objetos de tipos diferentes.

\subsection*{Escala de Planck, funcional de Barbero--Immirzi y mezcla CKM--CP}

La doble proyección precede, en el orden causal, a las construcciones
angulares y físicas. La forma normal de Smith de \(H_4\) produce un conúcleo
de orden dieciséis y la transferencia
\[
 \Sigma_H=\varphi\alpha_{\rm HMT}^{16}
\]
demuestra
\(\lfloor\log_{10}\Sigma_H\rfloor=-34\) y fija así la escala interna
\(10^{-34}\mathcal U_S\). La recta de acción contiene dos
coordenadas,
\[
 \hbar_{\rm pre}=H_5\,10^{-34}\mathcal U_S,\qquad
 \hbar_{\rm ret}=\eta_{\rm ret}\,10^{-34}\mathcal U_S.
\]
La periodicidad temporal de orden \(108\) proporciona una clausura autodual
de Planck. En la realización circular, con
\(\ell_0:=c_{\rm int}t_0\),
\[
 T_{108}=108t_0,
 \qquad
 R_{108}=\frac{54}{\pi}\ell_0,
 \qquad
 \theta_{108}^{\rm circ}=\left(\frac{54}{\pi}\right)^2,
\]
y el cuanto satisface
\[
 R_{108}=\bar\lambda_{\rm C}=r_{\rm g}=\ell_{\rm P}^{\rm circ}.
\]
La recta partida
\[
 Z_a(m)=\frac{m_{{\rm P},a}^{\circ}}mP_+
       +\frac m{m_{{\rm P},a}^{\circ}}P_-,
 \qquad N_{\rm sp}(Z_a(m))=1,
\]
intercambia Compton y gravedad bajo
\(m\mapsto(m_{{\rm P},a}^{\circ})^2/m\). Su punto fijo es la escala de
Planck. La sección electrónica constituye una realización genealógica
posterior sobre esa recta y no una identificación del electrón con un agujero
negro clásico.
La coordenada angular directa \(A\) y su coordenada angular conjugada
\(C^*\) construyen después los canales
\(q_\pm\). La inclusión hexada--octada produce el funcional de
Barbero--Immirzi, el operador de área y su espectro. La misma incidencia de
rango tres construye CKM, la fase CP y el invariante de Jarlskog. Estos
resultados pertenecen al cierre matemático de HMT; sus interpretaciones de
campo se retoman en MD.

La forma normal de Smith fija de manera explícita
\[
 \operatorname{SNF}(H_4)=\operatorname{diag}(1,2,2,4),
 \qquad |\operatorname{coker}H_4|=16,
\]
y la recta de acción distingue
\(\hbar_{\rm pre}\), \(\hbar_{\rm ret}\) y
\(h_{\rm ret}=2\pi\hbar_{\rm ret}\). La carta
\(\mathcal U_S\mapsto\mathrm{J\,s}\) publica después la unidad
metrológica; el orden \(-34\) pertenece a la sección interna que precede a
esa carta.

La inclusión combinatoria de una hexada en su octada asociada produce
\[
 6\binom62=90,\qquad
 8\binom62=120,
 \qquad
 \frac{90-120}{24\binom62}=-\frac1{12}.
\]
Sobre estos espacios, la cadena fundamental dodecafásica reúne doce sectores
y una única costura orientada de retorno. Su lector lineal produce
\(12S_{90}-S_{120}\) y, por tanto, los pesos \(12:-1\) del funcional
\(\Gamma_{6\to8}\). El coeficiente de \((1-q)^{-1}\), igual a \(1/24\),
la positividad y la minimalidad constituyen comprobaciones aritméticas
posteriores del par producido. La especialización de Holst
\(\gamma=\Gamma_{6\to8}\) se realiza en la frontera triádica mediante
\[
 \widehat A_{\rm HMT}
 =\gamma_{\rm HMT}a_0\sum_{e\in\partial A}N_e,
 \qquad
 \operatorname{Spec}(\widehat A_{\rm HMT})
 =\{n\gamma_{\rm HMT}a_0:0\le n\le72\}.
\]
El entrelazador colectivo satisface
\(J_{6\to8}D_{\rm inc}=D_{\rm area}J_{6\to8}\) en los subespacios
colectivos de incidencia. El operador de área, su espectro y las entropías
triádica, bicapa y modular constituyen así el cierre interno que realiza el
parámetro de Barbero--Immirzi. La transformación racional de rango tres
construye, en una rama hermana, los ángulos CKM, la fase CP y el invariante
de Jarlskog antes del contraste experimental.

\subsection*{Clausura genealógica del número y realización físico-espectral}

Las cuatro primeras Partes culminan donde una exposición convencional suele comenzar:
define el número como sección compatible y el valor como su proyección.
Mecánica Dimensional conserva esa definición y adopta un lector
físico-espectral. Primero
reconstruye el núcleo APP--TRIT--TPK en clave físico-espectral. APP se realiza
como toro bidimensional y distingue hojas aditiva, multiplicativa, areal,
volumétrica y torsional. TRIT recibe la semántica
\[
 \begin{aligned}
 +1&=\text{retorno, memoria y pasado},\\
 0&=\text{umbral y presente},\\
 -1&=\text{apertura y futuro compatible}.
 \end{aligned}
\]
TPK transporta la ruta, el registro y la holonomía. La periodicidad exacta
construye la dinámica finita; un cristal de tiempo físico añade Hamiltoniano,
subarmonía robusta y estabilidad. Las coordenadas angulares conjugadas
\(A\) y \(C^*\) realizan espectralmente la diferencia suma--producto.

La completación gnomónica adquiere aquí una función geométrica explícita. La
L porta la carta euclídea sobre la que los lectores inercial y gravitatorio
se proyectan a un cociente unitario. Esa carta constituye la realización del
principio local de equivalencia. El estado enriquecido anterior a la
proyección conserva separadas las hojas areal y volumétrica, y el Principio
de Ambivalencia fija su razón
\(\pi_{\rm HMT}/\varphi_{\rm HMT}^{2}\). Junto con las cartas de curvatura
negativa y positiva, se obtiene un atlas AdS--euclídeo--dS: la primera
realiza la correspondencia entre interior y borde, la segunda gobierna la
física local y la tercera satisface \(\ddot a/a=H^2>0\) y explica
geométricamente la expansión acelerada. Las tres cartas conviven porque son
realizaciones de una misma terna del TRIT, enlazadas por el transporte del
TPK.

Sobre esta gramática se define la partícula. Una partícula-ruta es una sección
persistente provista de orientación, hoja, ocupación, acción, propagador,
registro y anti-sección. La ley de masa se formula primero sobre cada
multisección \(F\):
\[
 \widehat{\mathbf M}^{\beta,r}_F
 =\mathbf m_\star^{\rm ret}\otimes
  R_{\rm act}^{\widehat K_F^r/2}
  \widehat{\mathcal A}^{(0)}_F
  R_{\rm act}^{\widehat K_F^r/2},
 \qquad r\in\{D,\Omega\},
\]
donde
\[
 \mathbf m_\star^{\rm ret}
 =\eta_{\rm ret}\,10^{-34}\mathcal U_S
   \frac{2\pi}{108t_0}\,c_{\rm int}^{-2}
\]
fija la sección absoluta de masa y
\(\widehat{\mathcal A}^{(0)}_F\) reúne firma, carácter y todas las torres de
\(\langle90,120\rangle\). El atlas completo es así el dominio de la ley; cada
especie selecciona después su realización dentro de la fibra correspondiente.

El electrón central es la primera reducción escalar, de rango uno, de esta
definición. La torsión global que interviene en su lector es
\[
 \Delta_4=\frac{\pi-e\log\pi}{270}
 \left(1+\frac{\pi}{729}\right).
\]
Su cierre bidireccional produce
\[
 \begin{aligned}
 R_{\rm act}&=\frac{H_5}{\eta_{\rm ret}},
 &\kappa_e&=80-54\alpha_{\rm HMT}+6\Delta_4,\\
 \mathfrak e_e^\beta&=\mathfrak e_e^{(0)}R_{\rm act}^{\kappa_e},\\
 \mathfrak e_e^\beta
 &=0.5109989506900008086082571648\ldots\ {\rm MeV},
 &m_e^\beta&=\mathfrak e_e^\beta/c^2.
 \end{aligned}
\]
La identidad
\(\Sigma_H=\varphi\alpha_{\rm HMT}^{16}\) aporta simultáneamente la
escala \(10^{-34}\mathcal U_S\) a las dos coordenadas de acción; el
cociente \(R_{\rm act}\) conserva el refinamiento relativo y cancela esa
escala común. Así, la escala interna de acción, la corrección bidireccional y la
carta energética en MeV ocupan niveles sucesivos de una misma derivación.
La primera cantidad es la energía de masa publicada; la segunda es la masa
obtenida por la carta explícita $E=mc^2$. Esta distinción impide identificar
una coordenada energética con una unidad de masa.
Las vacancias construyen polarización y corriente; la respuesta de las hojas
construye la permisividad y la permeabilidad del vacío. La carta SI publica
\[
 \begin{aligned}
 \mu_0^{\rm HMT}
 &=1.256637062121047789\ldots\times10^{-6}\,\mathrm{H\,m^{-1}},\\
 \varepsilon_0^{\rm HMT}
 &=8.854187812793002329\ldots\times10^{-12}\,\mathrm{F\,m^{-1}}.
 \end{aligned}
\]
El atlas sitúa después
los quarks y las restantes especies, definiendo color, sabor, carga, espín,
estadística y representación dentro de la misma genealogía.

En esta clasificación, el \emph{color} es la fibra ternaria que organiza las
tres realizaciones de una ruta quark y cuya composición neutra satisface
\(0+1+2=0\pmod3\). El \emph{sabor} es la coordenada que distingue, dentro
de una misma representación, las secciones transportadas por el operador de
mezcla. El \emph{espín} etiqueta la representación del recubrimiento del
grupo de rotaciones que porta la sección. La holonomía de la banda orientada
realiza sus cierres \(2\pi/4\pi\) y distingue el retorno de fase del retorno
de orientación. La estadística procede de la
graduación y de las realizaciones simétrica o exterior; fijadas CCR o CAR,
publica ocupación Bose--Einstein o Fermi--Dirac y, en el segundo caso, la
exclusión de Pauli. Una partícula virtual es una prolongación finita que
conserva ruta, registro y firma hasta un horizonte declarado; una sección
material prolongable conserva la compatibilidad global correspondiente.

\subsection*{Dominio operatorio de rutas, fibras y realizaciones físicas}

La construcción demuestra que la genealogía interna alcanza el dominio
físico completo. El dominio distingue
\(81\) posiciones, \(56\) familias, \(13\) multisecciones y \(324\) rutas
orientadas, junto con el sector nulo. Cada realización conserva, según su codominio,
posición o fibra, familia, ruta, orientación, palabra de fase, registro,
firma, carácter, operadores de refinamiento, carta y contraste. Los valores
externos se unen por identificador después de congelar la salida interna.

La cadena de construcción que gobierna esas filas es
\[
 \begin{aligned}
 \text{sección superviviente}&\to\text{ruta observable}
 \to\text{registro}\to\sigma\in\mathbb Z^5,\\
 \sigma&\to\chi_0\to\chi_{\rm full}\to\text{torres}
 \to\text{sección dimensional}.
 \end{aligned}
\]
Los dos lectores bidireccionales publican
\[
 K_D=\left(D_{27}+D_{36},D_1,\frac{D_4-D_3}{2}\right),
 \qquad
 K_\Omega=(\Omega,54,P_6),
\]
y actúan como operadores diagonales sobre cada fibra. En la celda central
ambos reducen a \((80,54,6)\); en las fibras no centrales sus espectros
conservan las distintas rutas antes de la selección de una realización
física. La comparación operatoria presenta primero la salida HMT, después la
carta y finalmente la referencia del Modelo Estándar o PDG con su unidad,
esquema, incertidumbre y residual.

La salida científica adopta el tipo que corresponde a cada sector. La celda
electrónica publica la masa beta escalar cerrada; fotón y gluones publican la
carta nula; las fibras no centrales conservan sus operadores y espectros.
El morfismo finito sabor--generación--color a ruta persistente ya está
construido, y las doce rutas nominales
\(e,\mu,\tau,u,d,s,c,b,t,W,Z,H\) publican sus firmas angulares y sus
caracteres escalares antes del contraste externo. En leptones y quarks, los
operadores \(K_D/K_\Omega\) refinan además la fibra; no son el requisito que
crea retroactivamente su escalarización. \(W\), \(Z\) y \(H\) conservan su
lector CT--108 nativo aunque no formen parte del atlas operatorial de trece
multisecciones. Los sectores topológicos publican fusión y trenzado; las
secciones virtuales publican su horizonte de prolongación. Esta
heterogeneidad preserva el codominio propio de cada clase y convierte la
tabla en una comparación completa de resultados tipados, valores y
fronteras de realización.

\subsection*{Cuántica, información, gravedad y cosmología}

La propagación de rutas construye la suma finita de caminos, la
representación Clifford--Dirac y la evolución de Schrödinger. Fijado el
Hamiltoniano del calendario interno, la evolución satisface
\[
 i\hbar_{\rm int}\dot\psi=H_{\rm clk}\psi,
\]
y su propagador finito admite la expansión
\[
 (U^R)_{fi}=
 \sum_{\gamma:i\to f,\,|\gamma|=R}
 \prod_{\ell\in\gamma}U_\ell.
\]
El producto conserva el orden de las transiciones dentro de una historia y
la suma acumula historias alternativas con extremos comunes. La graduación
del intercambio se realiza concretamente sobre
\(\mathbb T_{27}^{\,2}\times\{\mathrm N,\mathrm S,\mathrm E,\mathrm W\}\)
mediante
\[
 U_{\rm micro}=S D_J C_4,
 \qquad U_{12}^{\rm reg}=U_{\rm micro}^{12}.
\]
Este propagador es unitario; sus potencias enumeran exactamente las rutas
finitas y su transformada discreta produce \(27^2\) bloques de Fourier. La
construcción no se identifica por sí sola con una integral funcional
continua. La graduación fija el dato de intercambio; la representación del
grupo de permutaciones determina la completación simétrica o exterior. Dadas
después las relaciones CCR o CAR, el principio variacional y la maximización
de entropía producen las
distribuciones de Bose--Einstein y Fermi--Dirac, y la segunda realiza la
exclusión de Pauli. En la base Möbius, una vuelta visible de \(2\pi\) cambia
de hoja y acumula \(h/2\); el doble retorno visible de \(4\pi\) restaura la
orientación y acumula \(h\). Por radián de fase interna, el transporte
permanece igual a \(\hbar\).

La constante de Boltzmann aparece como transductor entre acción e información
ternaria. La construcción interna fija el producto
\(k_B\Theta_{\rm clk}\) y su función estructural; la definición del kelvin
del BIPM fija después la coordenada numérica de \(k_B\) en la carta SI:
\[
 k_B\Theta_{\rm clk}\log3=\hbar_{\rm int}\frac{2\pi}{108t_0},
 \qquad
 k_B^{\rm SI}=1.380649\times10^{-23}\ \mathrm{J\,K^{-1}}.
\]
Los ensambles realizan Gibbs y la condición KMS. El Landauer nulo corresponde
al transporte biyectivo de la fibra completa; preparación, proyección y
reinicio reciben balances separados cuando el acoplamiento descarta una
fibra. La lectura térmica construye primero el transductor interno y aplica
después la carta de unidades.

Las respuestas bicapa producen directamente impedancia, velocidad,
permeabilidad y permitividad mediante
\[
 \widehat Z=\sqrt{\widehat\mu/\widehat\varepsilon},
 \qquad
 \widehat c=(\widehat\mu\widehat\varepsilon)^{-1/2}.
\]
La misma precedencia causal gobierna la lectura electromagnética: el operador
constitutivo interno antecede a la realización de
\(\varepsilon_0,\mu_0,Z_0\) y \(c\) en el Sistema Internacional.

La rama gravitatoria construye primero, desde las salidas del estado
enriquecido, la escala de acción, el selector gravitatorio y la sección
\(G_a\). La formulación variacional de Palatini--Cartan--Holst recibe después
esas salidas, junto con \(\gamma_{\rm HMT}\), y admite la transformada
hamiltoniana restringida
\[
 H_{\rm tot}=\int_\Sigma
 (\lambda^i G_i+N^a V_a+N\mathcal H).
\]
La equivalencia posterior entre acción covariante, sistema hamiltoniano restringido y
evolución Einstein--Schrödinger conserva co-tétrada, conexión, curvatura y
torsión dentro del mismo esquema. El sector colectivo libre de torsión porta
la excitación gravitatoria; la torsión local inducida por la densidad de espín
aporta el canal material. El tipo, el transductor y la carta de \(G\) anteceden
a esta realización dinámica y no se seleccionan mediante ella.

Las ramas solenoidal y gauge reciben realizaciones del operador finito común
\[
 \mathcal L_{\rm cyc}^{\rm TPK}=2I-C-C^{-1}=3(I-P_0)=3I-J,
\]
\[
 \mathfrak R_{\rm sol}(\mathcal L_{\rm cyc}^{\rm TPK})
 =\mathcal L_{\rm H}^{(0)},\qquad
 \mathfrak R_{\rm gau}(\mathcal L_{\rm cyc}^{\rm TPK})
 =\mathcal L_{\rm H}^{(1)}.
\]
La rama hidrodinámica lo transporta mediante el entrelazador
Galerkin--registro, con coercividad y residual declarados; la rama gauge lo
transporta mediante el funcional gauge y el operador
transferencia--brecha. La demostración finita y los puentes a las ecuaciones
continuas se exponen con sus hipótesis diferenciadas. Los marcos
autoinerciales, la factorización machiana y la cosmología de hojas se
desarrollan después de esta construcción local y colectiva, donde adquieren
dominio, lector e interpretación física tipada.

\subsection*{Método de demostración y contraste}

Cada resultado se presenta mediante objeto, dominio, mapa, derivación,
premisas, certificado y alcance. Los cálculos reproducibles acompañan a la
demostración y no la sustituyen. Las cartas metrológicas se aplican después
de construir el escalar o el operador interno. Las comparaciones CODATA y PDG
se abren después del sellado de la salida HMT\@. Este orden preserva la
dirección causal del programa: la estructura genera la cifra y la magnitud;
la referencia externa mide el acuerdo alcanzado por HMT\@.

La obra se organiza en cinco Partes causalmente ordenadas. Las cuatro
primeras construyen el núcleo formal, el continuo discreto, la genealogía de
las constantes y las cinco cadenas de especialización; la quinta realiza esa arquitectura
en magnitudes, partículas, campos y geometría física. Esta secuencia expresa
la dependencia matemática entre estructura, evaluación, clausura y
realización dimensional.

\subsection*{Composición del tratado y alcance de las afirmaciones}

La Parte I construye APP desde la célula nonádica, presenta TRIT como
levantamiento orientado y define el Topological Phase Kernel con sus fibras,
transportes, memoria, reloj, conexión nonádica conjunta y prolongación
compatible. La línea del horizonte abre esta construcción mediante las
nociones de escala, frontera y medida.

La Parte II desarrolla la estructura discreta del continuo. La doble
proyección separa la evaluación arquimediana de la publicación incidencial.
Desde Paley--Witt se distinguen la rama de diseños Steiner--Golay y la rama
reticular Niemeier--Leech--\(V^\natural\)--Monstruo--Moonshine. Dentro de esa misma estructura
aparecen la raíz interna de \(-1\), las tres realizaciones cuadráticas y las
subestructuras espectral--polar, solenoidal, gauge, cohomológica y
determinantal, tipadas por sus propios mapas y terminales.

La Parte III establece la genealogía de las constantes fundamentales: la
publicación coinductiva correlacionada de
\(\pi,e,\varphi,\alpha_{\rm HMT}\), los lectores específicos y las dos vías
internas de \(\alpha_{\rm HMT}\), \(-1/12\), la escala de acción, Planck,
Boltzmann, el parámetro de Barbero--Immirzi, CKM--CP, las constitutivas del
vacío, la constante gravitatoria y el atlas de constantes matemáticas y
físicas.

\Needspace{12\baselineskip}
La Parte IV aplica las cinco subestructuras construidas en la Parte II en
capítulos de idéntico rango editorial. Cada capítulo expone el resultado
interno, el comparador especializado, sus premisas y el criterio exacto de
publicación. La conclusión en la formulación convencional correspondiente se
deduce únicamente cuando la identidad comparadora queda verificada en todo su
dominio.

La Parte V realiza APP--TRIT--TPK mediante los lectores físicos: toro y hojas,
semántica temporal, curvaturas, canales y holonomía. Sobre esa gramática
define la partícula-ruta, el electrón bidireccional, las familias de masa,
las constitutivas, la estadística, la información, la mezcla, la gravedad
variacional, la autoinercia, la cosmología de hojas y la interfaz con teoría
M.

Los apéndices reúnen demostraciones auxiliares, censos completos, semillas,
cadenas decimales, tablas extensas, derivaciones auxiliares y cálculos
reproducibles. El cuerpo principal conserva la secuencia definición,
construcción, teorema, interpretación y resultado. Cada afirmación distingue
su dominio matemático, su estructura de partida, su carta de realización y su
contraste externo. Esta separación permite presentar juntos los cierres
exactos, las realizaciones físicas y los programas de extensión, conservando
para cada uno el grado de demostración que le corresponde.
